\documentclass[11pt,a4paper]{article}
\usepackage[margin=2.6cm]{geometry}
\usepackage[utf8]{inputenc}
\usepackage[T1]{fontenc}
\usepackage{lmodern}
\usepackage{microtype}
\usepackage{amsmath,amssymb,amsthm,mathtools}
\usepackage{enumitem}
\usepackage{array}
\usepackage{hyperref}
\usepackage{url}
\hypersetup{colorlinks=true,linkcolor=blue,citecolor=blue,urlcolor=blue}

\newtheorem{theorem}{Theorem}[section]
\newtheorem{lemma}[theorem]{Lemma}
\newtheorem{proposition}[theorem]{Proposition}
\newtheorem{corollary}[theorem]{Corollary}
\newtheorem{definition}[theorem]{Definition}
\newtheorem{remark}[theorem]{Remark}
\newtheorem{example}[theorem]{Example}
\newtheorem{observation}[theorem]{Observation}
\theoremstyle{remark}
\newtheorem*{remark*}{Remark}

\newcommand{\CS}{\mathrm{CS}}
\newcommand{\Cmcf}[2]{\mathcal{C}^{\mathrm{mcf}}_{#1,#2}}
\newcommand{\Ccfg}[1]{\mathcal{C}^{\mathrm{cfg}}_{#1}}
\newcommand{\Yclass}[1]{\mathcal{Y}_{#1}}
\newcommand{\Yord}[1]{\mathcal{Y}^{\mathrm{ord}}_{#1}}
\newcommand{\Yyos}[2]{\mathcal{Y}^{\mathrm{Yos}}_{#1,#2}}
\newcommand{\Dord}{\mathcal{D}}
\newcommand{\D}{\mathcal{D}}
\newcommand{\out}{\mathrm{out}}
\newcommand{\ev}{\mathrm{ev}}
\newcommand{\id}{\mathrm{id}}
\newcommand{\st}{\mathrm{st}}
\newcommand{\tuple}[1]{\vec{#1}}
\newcommand{\blank}{\square}
\newcommand{\eps}{\varepsilon}
\newcommand{\poly}{\mathrm{poly}}
\newcommand{\yield}{\mathrm{yield}}
\newcommand{\type}{\mathrm{type}}
\newcommand{\rank}{\mathrm{rank}}
\newcommand{\Sp}{\mathrm{Sp}}
\newcommand{\Lex}{\mathrm{Lex}}

\providecommand{\keywords}[1]{\par\smallskip\noindent\textbf{Keywords. }#1\par}



\title{Positive-Data Learning of Bounded-Fan-Out Linear MCFGs with Fixed Finite-Monoid Advice}
\author{Takayuki Kuriyama\\Independent Researcher, Tokyo, Japan}
\date{}


\begin{document}
\setlength{\abovedisplayskip}{6pt plus 2pt minus 3pt}
\setlength{\belowdisplayskip}{6pt plus 2pt minus 3pt}
\setlength{\abovedisplayshortskip}{3pt plus 1pt minus 2pt}
\setlength{\belowdisplayshortskip}{3pt plus 1pt minus 2pt}
\maketitle



\begin{abstract}
We study positive-data learning of a semantic subclass of bounded-fan-out
linear multiple context-free languages under fixed finite-state side
information.  The side information is an explicit finite monoid homomorphism
\(h:\Sigma^*\to M\), supplied before learning starts and not inferred from the
positive text.

For reduced working binary linear nondeleting MCFG presentations, we define
\((f,h)\)-tuple substitutability using named sentence contexts and
componentwise \(h\)-types.  For every fixed fan-out bound \(f\) and fixed
explicit morphism \(h\), we construct a canonical set-driven learner.  The
proof uses an output-type refinement of a witnessing presentation together
with concrete shared sentence contexts exposed by a finite
presentation-relative characteristic sample.  The refinement records
componentwise output types, while the sample supplies the information needed
to recover component placement and to guard substitution.  If the target admits
the required working presentation and satisfies the semantic
\((f,h)\)-substitutability promise, then every positive sample containing the
characteristic sample yields a hypothesis generating exactly the target
language.  Consequently, each fixed observation fiber is identifiable in the
limit from positive data.  For fixed \((f,h)\), the hypothesis associated with
a finite sample \(K\) is constructible in time \(\|K\|_+^{O(f)}\), including
output size.

We also determine a boundary for the observation advice.  All observations
whose monoids have a fixed size bound can be compiled into one finite product
morphism and remain identifiable, whereas the unbounded union over all finite
observations is not identifiable from positive data.  Finally, an infinite
member-kernel criterion excludes languages, including the copy language, from
every finite-observation fiber.  Thus fixed finite-monoid advice isolates a
proper learnable region inside the multiple context-free languages rather than
a reconstruction of the full class.
\end{abstract}

\keywords{grammatical inference, multiple context-free grammars, positive data,
identification in the limit, finite monoids, tuple substitutability}



\section{Introduction}
\label{sec:introduction}

Gold's theorem shows that broad language classes cannot be identified from
positive data alone \cite{Gold1967}.  A successful way to recover positive-data
learnability is to impose distributional restrictions: observed substrings, or
more generally observed tuples, may be safely substituted when their accepting
contexts agree in a sufficiently controlled way.  This viewpoint underlies the
substitutable context-free languages of Clark and Eyraud \cite{ClarkEyraud2007}
and the multidimensional substitutability results of Yoshinaka for mildly
context-sensitive languages \cite{Yoshinaka2008,Yoshinaka2009,Yoshinaka2011}, as
well as later work on distributional learning of parallel and multiple
context-free grammars \cite{ClarkYoshinaka2014,ClarkYoshinaka2016}.  Related
positive-data learnability phenomena also appear in restricted categorial
grammars \cite{Kanazawa1998}.  For the mildly context-sensitive grammar
background motivating MCFGs and related formalisms, see
\cite{Shieber1985,VijayShankerWeirJoshi1987,Weir1988,Kallmeyer2010}.

This paper studies a typed version of that idea for bounded-fan-out working
multiple context-free grammars (MCFGs).  It can be read as a mildly
context-sensitive analogue of fixed finite-monoid typed distributional learning
for context-free languages \cite{kuriyamaCFG}.  The type information is supplied
by a fixed explicit finite monoid homomorphism \(h:\Sigma^*\to M\).  The
homomorphism is not inferred from the text.  It is observation advice: a finite
interface that may record, for instance, a regular envelope or a coarse
finite-state annotation.  We use ``advice'' only in the learning-theoretic sense
of fixed side information supplied before the text; no nonuniform
complexity-theoretic advice class is involved.  The learner is therefore
advice-relative.  It learns the target language from positive examples, but the
observation interface itself is fixed in advance.  All results are stated for languages admitting reduced
working binary linear nondeleting MCFG presentations of bounded fan-out.  No
normal-form theorem for arbitrary MCFGs is used or claimed.

\paragraph{What is learned, and what is not.}
The learner receives positive examples of the target language and the fixed
finite morphism \(h\).  It does not receive a target grammar, a derivation
oracle, negative examples, membership queries, or the semantic promise as
decidable input.  The theorem says that if the unknown target belongs to the
fixed fiber \(\Cmcf{f}{h}\), then a finite positive sample is sufficient for the
canonical set-driven learner to stabilize to the exact target language.  It does
not say that \(h\) can be discovered from positive data, nor that arbitrary
MCFGs are identifiable from positive data.

The target condition is called \((f,h)\)-tuple substitutability.  Since MCFG
nonterminals derive tuples rather than single strings, the relevant contexts are
named sentence contexts with several holes.  Two tuples of arity at most \(f\)
may be identified only when they have the same componentwise \(h\)-type and
share at least one accepting sentence context; the semantic promise then says
that all accepting contexts coincide.  This is a semantic promise on the target
language, not a decidable syntactic promise on an input grammar.

The main contributions are as follows.
\begin{enumerate}[label=(\roman*),leftmargin=*]
\item We define the fixed observation fibers \(\Cmcf{f}{h}\) for languages
admitting reduced working binary linear nondeleting MCFG presentations and
satisfying semantic \((f,h)\)-tuple substitutability.

\item For each fixed \(f\) and fixed explicit finite morphism \(h\), we construct
a canonical positive-data learner and prove exact language reconstruction from
a finite presentation-relative characteristic sample.  The learner stores only
observed tuple values; component placement is recovered from concrete witnessed
segmentations, and substitution is guarded by actually shared sentence
contexts.

\item We separate construction time from data size.  The hypothesis associated
with a given finite sample is constructible in \(\|K\|_+^{O(f)}\) time,
including output size.  Polynomial-size characteristic samples follow for
subclasses with a polynomial exposure bound, but no such bound is asserted for
the full target class.

\item We prove that finite observation is genuine advice: the union of all
observations with a fixed monoid-size bound compiles into one product morphism,
whereas the unbounded union over all finite morphisms is not identifiable from
positive data.

\item We give a presentation-independent member-kernel exclusion criterion and
apply it to a context-free slope union and to the copy language.  Thus the
fixed-observation fibers form a proper advice-relative fragment of the MCFLs
rather than a reformulation of the full hierarchy.
\end{enumerate}

Section~\ref{sec:main-results} states the theorem package.
Section~\ref{sec:prelim} gives the learning and grammar-theoretic definitions,
and Section~\ref{sec:running-examples} records standard synchronization
examples.  Sections~\ref{sec:learner} and \ref{sec:exact} define the learner
and prove exact identification.  Section~\ref{sec:poly-time} separates the
fixed-parameter hypothesis-construction bound from the exposure condition
needed for polynomial data.  Section~\ref{sec:fixed-observation} proves the
bounded/unbounded observation-advice boundary and the member-kernel exclusions.
Section~\ref{sec:comparison} compares the result with untyped distributional
conditions.  Appendix~\ref{app:semantic-promise} records the non-effectivity of
the semantic promise.


\section{Main Results}
\label{sec:main-results}

This section states the theorem package before the technical development.  The
proofs are given in Sections~\ref{sec:learner}--\ref{sec:fixed-observation}.
Throughout, the working-presentation assumptions are exactly those of
Definition~\ref{def:working-mcfg}; no normalization theorem for arbitrary MCFGs
is being invoked.

\begin{theorem}[Exact reconstruction, preview]
\label{thm:main-exact-preview}
Fix a fan-out bound \(f\) and an explicit finite homomorphism
\(h:\Sigma^*\to M\).  Let \(G\) be a reduced working binary linear
nondeleting MCFG of fan-out at most \(f\), and assume that
\(L=L(G)\) is \((f,h)\)-tuple-substitutable.  Then a finite
presentation-relative characteristic sample \(S_G\subseteq L\) exists such
that, for every finite \(K\) with \(S_G\subseteq K\subseteq L\), the canonical
learner returns a grammar whose language is exactly \(L\).
\end{theorem}

This statement is proved as Theorem~\ref{thm:exact-reconstruction};
Corollary~\ref{cor:gold} gives identification in the limit.

\begin{theorem}[Hypothesis construction, preview]
\label{thm:main-poly-preview}
For fixed \(f\) and fixed explicit \(h\), the canonical hypothesis associated
with a finite sample \(K\) is constructible in time
\(\|K\|_+^{O(f)}\), including its output size.  This is a construction bound
from the sample at hand; it is not by itself a polynomial-data theorem for the
full target class.
\end{theorem}

The construction bound is proved in Theorem~\ref{thm:poly}, and the precise
polynomial-data boundary is recorded in
Proposition~\ref{prop:cs-exposure-bound}.

\begin{theorem}[Observation-advice boundary, preview]
\label{thm:main-advice-preview}
For every fixed alphabet \(\Sigma\), fan-out bound \(f\), and monoid-size bound
\(k\), the union of all fibers defined by morphisms into monoids of size at
most \(k\) is identifiable from positive data.  In contrast, the union over
all finite observation morphisms, with no size bound, is not identifiable from
positive data.
\end{theorem}

The two parts are Theorems~\ref{thm:bounded-union-identifiable} and
\ref{thm:no-advice-nonidentifiability}.

\begin{theorem}[Finite-observation exclusion, preview]
\label{thm:main-exclusion-preview}
If the arity-one accepting sentence-context distributions of members of a
language have infinite index, then the language belongs to no fixed
finite-observation fiber.  In particular, the copy language is outside every
such fiber.
\end{theorem}

This is Theorem~\ref{thm:member-kernel-exclusion} together with
Proposition~\ref{prop:copy-language-outside-every-finite-observation}.

\section{Preliminaries and the Target Class}
\label{sec:prelim}

\subsection{Positive-data learning}

\begin{definition}[Text and identification]
A \emph{text} for \(L\subseteq\Sigma^*\) is an infinite sequence
\((w_1,w_2,\ldots)\) of elements of \(L\) in which every element of \(L\)
appears at least once.  A learner is a computable function mapping each finite
prefix to a hypothesis grammar.  A learner identifies a class \(\mathcal C\) in
the limit from positive data if for every \(L\in\mathcal C\) and every text for
\(L\), there exists \(n_0\) such that all hypotheses produced after stage
\(n_0\) have language \(L\).
\end{definition}

The learner below is order-independent, so a finite text prefix is identified
with its finite set \(K\subseteq L\).

\begin{definition}[Characteristic sample]
For a learner \(\mathcal A\) and a target language \(L\), a finite set
\(S\subseteq L\) is a \emph{characteristic sample} if
\(L(\mathcal A(K))=L\) for every finite \(K\) with \(S\subseteq K\subseteq L\).
\end{definition}

\begin{definition}[Positive sample size]
For finite \(K\subseteq\Sigma^*\), define \(\|K\|_+:=\sum_{w\in K}\max(1,|w|)\).
\end{definition}

The factor \(\max(1,|w|)\) only prevents degenerate zero-size samples when an
empty word is allowed.  In the working nondeleting setting below, generated
nonempty examples are the main case.

\begin{definition}[Polynomial time and data]
\label{def:poly-time-data}
Following the characteristic-sample viewpoint of de~la~Higuera
\cite{delaHiguera1997,delaHiguera2010}, a presentation class is said here to be
identified in polynomial time and data, for fixed external parameters, if there
are polynomials \(p\) and \(q\) such that every target with a witnessing
presentation \(G\) has a characteristic sample of positive size at most
\(p(|G|)\), and the learner's hypothesis on every finite sample \(K\) is
constructible in time at most \(q(\|K\|_+)\), including output size.  The
polynomial-data part is presentation-relative: the bound is measured against a
chosen witnessing presentation, not against a canonical minimal description of
the language.
\end{definition}

\begin{remark}[Slicewise, not uniform, polynomiality]
\label{rem:slicewise-poly}
Throughout, ``polynomial time'' and ``polynomial data'' are understood
slicewise in the external parameters: for each fixed \(f\) and fixed \(h\), the relevant bound is a polynomial
whose degree may depend on those parameters.  Concretely, the construction time
is \(\|K\|_+^{O(f)}\).  This is the fixed-parameter reading of the
de~la~Higuera polynomial criterion under fixed external parameters, not a claim
of a single polynomial uniform in \(f\).  No uniform-in-\(f\) bound is asserted or
used.
\end{remark}

Gold identification \cite{Gold1967} is the limiting correctness notion used in
this paper; characteristic samples are the finite sufficient sets that make the
set-driven learner stabilize.  The terminology ``polynomial time and data'' is
therefore not an additional learning model, but the de~la~Higuera-style
polynomial strengthening of positive-data identification.  Angluin's exact
learning model \cite{Angluin1987} is cited only as a related landmark in
algorithmic learning; the learner in this paper receives no queries and no
counterexamples.

\begin{lemma}[Finite sufficient samples imply identification]
\label{lem:sufficient-to-gold}
Let \(\mathcal A\) be a computable order-independent learner.  Suppose that,
for each \(L\in\mathcal C\), there is a finite \(S_L\subseteq L\) such that
\(S_L\subseteq K\subseteq L\) implies \(L(\mathcal A(K))=L\).  Then
\(\mathcal A\) identifies \(\mathcal C\) in the limit from positive data.
\end{lemma}

\begin{proof}
In every text for \(L\), the finitely many words of \(S_L\) have all appeared
after some stage \(n_0\).  Every later prefix sample is contained in \(L\) and
contains \(S_L\), so the hypothesis language is \(L\) from that stage on.
\end{proof}

\subsection{Finite monoid typing}

\begin{definition}[Explicit finite homomorphism]
A monoid homomorphism \(h\colon\Sigma^*\to M\) into a finite monoid \(M\) is
\emph{explicit} if the multiplication table of \(M\) and the values \(h(a)\) for
\(a\in\Sigma\) are given.  For a tuple \(\tuple w=(w_1,\ldots,w_d)\), write
\(h^{(d)}(\tuple w):=(h(w_1),\ldots,h(w_d))\in M^d\).
\end{definition}

\begin{definition}[Refinement of finite observation morphisms]
\label{def:h-refinement}
Let \(h\colon\Sigma^*\to M\) and \(h'\colon\Sigma^*\to M'\) be explicit finite
homomorphisms.  We say that \(h'\) \emph{refines} \(h\), and write
\(h\preceq h'\), if there is a monoid homomorphism
\(\pi\colon M'\to M\) such that \(h=\pi\circ h'\).
\end{definition}

\begin{proposition}[Monotonicity under refinement of the observation morphism]
\label{prop:h-refinement-monotonicity}
If \(h\preceq h'\) and \(L\) is \((f,h)\)-tuple-substitutable, then \(L\) is
\((f,h')\)-tuple-substitutable.  Consequently, for every pair of explicit finite
homomorphisms with \(h\preceq h'\),
\[
  \Cmcf{f}{h}\subseteq \Cmcf{f}{h'} .
\]
\end{proposition}

\begin{proof}
Let \(\pi\colon M'\to M\) satisfy \(h=\pi\circ h'\).  Suppose that
\(h'^{(d)}(\tuple x)=h'^{(d)}(\tuple y)\) and that \(\tuple x,\tuple y\) share an
accepting arity-\(d\) sentence context.  Applying \(\pi\) componentwise gives
\(h^{(d)}(\tuple x)=h^{(d)}(\tuple y)\).  The \((f,h)\)-tuple-substitutability of
\(L\) therefore gives \(\D_L^{(d)}(\tuple x)=\D_L^{(d)}(\tuple y)\).  This is
exactly the \((f,h')\)-condition.  The inclusion of classes follows because the
same witnessing working presentation may be used.
\end{proof}

Finite monoids are used here only as fixed external typing devices; see
Pin~\cite{Pin1986} for background.  Typical examples include transition
monoids of regular control languages, syntactic morphisms of regular
approximations, and finite-state annotations supplied by an external parser or
formalism.  The learner below assumes such an explicit morphism as part of its
fixed parameters.  In particular, when \(h\) is the transition morphism of a
regular envelope, the learner is being given that regular envelope in finite
form.  This is useful but nontrivial advice; the theorems below are relative to
that advice and do not solve the problem of discovering it from positive data.

\subsection{Working MCFG presentations}

\begin{definition}[Working binary linear nondeleting MCFG]
\label{def:working-mcfg}
A \emph{working binary linear nondeleting multiple context-free grammar} is a
tuple \(G=(V,\Sigma,P,S,\mu)\), where \(V\) is a finite nonterminal set,
\(\Sigma\) is a finite terminal alphabet, \(S\in V\) is the start symbol, and
\(\mu\colon V\to\mathbb N_{>0}\) is the fan-out map with \(\mu(S)=1\).  Rules
have one of the following forms.
\begin{enumerate}[label=(\roman*),leftmargin=*]
\item A start rule \(S\to A\), where \(A\ne S\) and \(\mu(A)=1\).
\item A terminal rule \(A\to(a)\), where \(A\ne S\), \(a\in\Sigma\), and
\(\mu(A)=1\).
\item A binary rule \(\rho\colon A\to(\alpha_1,\ldots,\alpha_e)(B,C)\), where
\(A\ne S\), \(e=\mu(A)\), \(\mu(B)=d_B\), \(\mu(C)=d_C\), each \(\alpha_i\) is
a word over \(\Sigma\cup\{x^1,\ldots,x^{d_B},y^1,\ldots,y^{d_C}\}\), and each
variable \(x^1,\ldots,x^{d_B},y^1,\ldots,y^{d_C}\) occurs exactly once in the
whole tuple \((\alpha_1,\ldots,\alpha_e)\).
\end{enumerate}
The presentation is \emph{start-separated}: the only rules with left-hand side
\(S\) are start rules, and \(S\) never occurs as a child on the right-hand side
of a binary rule.
\end{definition}

Consider a binary rule
\(\rho\colon A\to(\alpha_1,\ldots,\alpha_e)(B,C)\), where
\(e=\mu(A)\), \(d_B=\mu(B)\), and \(d_C=\mu(C)\).
Define the variable sets associated with this rule by
\(X_B:=\{x^1,\ldots,x^{d_B}\}\) and
\(Y_C:=\{y^1,\ldots,y^{d_C}\}\), and put
\(\Gamma_\rho:=\Sigma\uplus X_B\uplus Y_C\).
Each word \(\alpha_\ell\in\Gamma_\rho^*\), for
\(1\le \ell\le e\), is called a \emph{template component} of \(\rho\), and
\(\boldsymbol{\alpha}_\rho:=(\alpha_1,\ldots,\alpha_e)
\in(\Gamma_\rho^*)^e\) is called the \emph{template tuple} of \(\rho\).

For tuples
\(\tuple u=(u_1,\ldots,u_{d_B})\in(\Sigma^*)^{d_B}\) and
\(\tuple v=(v_1,\ldots,v_{d_C})\in(\Sigma^*)^{d_C}\), define
\(\sigma_{\tuple u,\tuple v}\colon\Gamma_\rho\to\Sigma^*\) by
\(\sigma_{\tuple u,\tuple v}(a):=a\) for \(a\in\Sigma\),
\(\sigma_{\tuple u,\tuple v}(x^i):=u_i\) for \(1\le i\le d_B\), and
\(\sigma_{\tuple u,\tuple v}(y^j):=v_j\) for \(1\le j\le d_C\).
Let \(\widehat{\sigma}_{\tuple u,\tuple v}\colon
\Gamma_\rho^*\to\Sigma^*\) be the unique extension of this map to a monoid
morphism.  The result of
applying \(\rho\) to \(\tuple u\) and \(\tuple v\) is defined by
\[
  \rho(\tuple u,\tuple v)
  :=
  \bigl(
    \widehat{\sigma}_{\tuple u,\tuple v}(\alpha_1),
    \ldots,
    \widehat{\sigma}_{\tuple u,\tuple v}(\alpha_e)
  \bigr)
  \in(\Sigma^*)^e.
\]
This is the simultaneous-substitution semantics of a binary rule.

For each nonterminal \(A\in V\), its tuple language
\(L_A(G)\subseteq(\Sigma^*)^{\mu(A)}\) is the component of the least family
\((L_A(G))_{A\in V}\), ordered by
componentwise inclusion, satisfying the following closure conditions.
\begin{enumerate}[label=(\roman*),leftmargin=*]
\item If \(A\to(a)\) is a terminal rule, then \((a)\in L_A(G)\).

\item If \(\rho\colon A\to(\alpha_1,\ldots,\alpha_e)(B,C)\) is a
      binary rule, \(\tuple u\in L_B(G)\), and
      \(\tuple v\in L_C(G)\), then
      \(\rho(\tuple u,\tuple v)\in L_A(G)\).

\item If \(S\to A\) is a start rule and \(\tuple w\in L_A(G)\), then
      \(\tuple w\in L_S(G)\).
\end{enumerate}
The string language generated by \(G\) is
\(L(G):=\{w\in\Sigma^*\mid (w)\in L_S(G)\}\).
Since \(\mu(S)=1\), we identify \(L_S(G)\) with \(L(G)\) when no confusion
can arise.

A nonterminal \(A\) is \emph{productive} if \(L_A(G)\ne\emptyset\).
The nonterminal-dependency graph of \(G\) has an edge \(S\to A\) for every
start rule \(S\to A\), and edges \(A\to B\) and \(A\to C\) for every binary
rule \(A\to(\alpha_1,\ldots,\alpha_e)(B,C)\).
A nonterminal \(A\) is \emph{reachable} if there is a directed path from
\(S\) to \(A\) in this graph.  The grammar \(G\) is \emph{reduced} if every
nonterminal is both reachable and productive.

\begin{remark}[Scope and non-claims of the working form]
\label{rem:working-form-scope}
Definition~\ref{def:working-mcfg} fixes the presentation model used by the
learner and by the characteristic-sample proof.  Every such presentation is a
standard MCFG presentation in the sense of Seki et al.~\cite{SekiEtAl1991}, but
we do not claim here that every standard MCFG can be effectively transformed
into this exact working form while preserving all parameters relevant to the
learning theorem.

The start-separated convention is harmless and can be enforced by introducing a
fresh start symbol.  By contrast, the binary, linear, and nondeleting
requirements are genuine presentation assumptions of this paper.  The theorem
therefore applies to languages admitting a reduced presentation in this working
form and satisfying the semantic \((f,h)\)-tuple-substitutability promise.

This restriction is intentional.  It keeps the observed binary witnesses in the
finite sample aligned with the binary rules reconstructed by the learner.  More
general ranks, deleting rules, and explicit \(\eps\)-generation would require
additional bookkeeping and are not part of the present claim.
\end{remark}

\subsection{Sentence contexts and tuple distributions}

\begin{definition}[Sentence context]
\label{def:sentence-context}
A \emph{sentence context of arity \(d\)} is a word of the form
\[
  E=u_0\blank_{\sigma(1)}u_1\cdots\blank_{\sigma(d)}u_d,
\]
where \(\sigma\) is a permutation of \(\{1,\ldots,d\}\), each
\(u_i\in\Sigma^*\), and each named hole
\(\blank_i\) occurs exactly once.  For \(\tuple w=(w_1,\ldots,w_d)\), write
\(E[\tuple w]\) for the string obtained by substituting \(w_i\) for
\(\blank_i\).
\end{definition}

Template components in later MCFG rules may be empty.  When empty tuple
components are observed inside a sample word, we represent them by cut positions
with a fixed local tie order; this is only bookkeeping for enumerating concrete
occurrences.

\begin{definition}[Tuple distribution]
\label{def:tuple-distribution}
Let \(d\ge1\).  For \(L\subseteq\Sigma^*\) and
\(\tuple x\in(\Sigma^*)^d\), define
\[
  \D_L^{(d)}(\tuple x)
  :=
  \left\{
    E
    \;\middle|\;
    E\text{ is an arity-}d\text{ sentence context and }
    E[\tuple x]\in L
  \right\}.
\]
When the arity is clear, write \(\D_L(\tuple x)\).
\end{definition}

\begin{definition}[\((f,h)\)-tuple substitutability]
\label{def:tuple-substitutable}
Let \(f\ge1\) and let \(h\colon\Sigma^*\to M\) be an explicit finite monoid
homomorphism.  A language \(L\subseteq\Sigma^*\) is
\emph{\((f,h)\)-tuple-substitutable} if, for every \(1\le d\le f\) and all
\(\tuple x,\tuple y\in(\Sigma^*)^d\), the implication
\(h^{(d)}(\tuple x)=h^{(d)}(\tuple y)\) and
\(\D_L^{(d)}(\tuple x)\cap\D_L^{(d)}(\tuple y)\ne\emptyset\) imply
\(\D_L^{(d)}(\tuple x)=\D_L^{(d)}(\tuple y)\).
\end{definition}

For \(d=1\), sentence contexts are ordinary two-sided contexts
\(u\blank_1v\), so Definition~\ref{def:tuple-substitutable} specializes to the
usual fixed-\(h\) two-sided substitutability condition.

\begin{definition}[The target class]
\label{def:class-cmcf}
Fix \(f\ge1\) and an explicit finite monoid homomorphism \(h\).  The class
\(\Cmcf{f}{h}\) consists of all languages \(L\subseteq\Sigma^*\) such that
\(L=L(G)\) for some reduced working binary linear nondeleting MCFG \(G\) with
all fan-outs at most \(f\), and \(L\) is \((f,h)\)-tuple-substitutable.
\end{definition}

\begin{remark}[Semantic nature of the target class]
\label{rem:semantic-class}
For a language \(L\) to belong to \(\Cmcf{f}{h}\), the following two
conditions must hold.
\begin{enumerate}[label=(\roman*),leftmargin=*]
\item There is a reduced working MCFG \(G\) of fan-out at most \(f\) such that
      \(L=L(G)\).
\item The language \(L(G)\) is \((f,h)\)-tuple-substitutable.
\end{enumerate}
Condition~(i) concerns a concrete grammar presentation.  By contrast,
condition~(ii) concerns the sentence-context distributions of the entire
language generated by that presentation, rather than the syntactic form of its
rules.  In this sense, \(\Cmcf{f}{h}\) is a \emph{semantic target class}.

We do not claim that, given a working MCFG \(G\), one can always decide whether
\(L(G)\) is \((f,h)\)-tuple-substitutable.  Failure of tuple substitutability is,
however, semi-decidable, because it has a finite witness.

More precisely, failure is witnessed by some \(1\le d\le f\), tuples
\(\tuple x,\tuple y\in(\Sigma^*)^d\), and arity-\(d\) sentence contexts
\(E,F\) such that
\[
\begin{gathered}
  h^{(d)}(\tuple x)=h^{(d)}(\tuple y),\\
  E[\tuple x],E[\tuple y],F[\tuple x]\in L(G),
  \qquad F[\tuple y]\notin L(G).
\end{gathered}
\]
The context \(E\) shows that \(\tuple x\) and \(\tuple y\) share at least one
accepting sentence context.  The context \(F\), on the other hand, shows that
\(F\in\D_{L(G)}^{(d)}(\tuple x)\setminus
\D_{L(G)}^{(d)}(\tuple y)\).  Consequently,
\(\D_{L(G)}^{(d)}(\tuple x)\ne
\D_{L(G)}^{(d)}(\tuple y)\), which violates Definition~\ref{def:tuple-substitutable}.

Finite tuples and finite sentence contexts can be effectively enumerated, for
example in increasing length-lexicographic order.  Since MCFG membership is
decidable, each candidate can be checked in finite time, including the required
nonmembership test.  Therefore, if a failure witness exists, an exhaustive
search eventually finds one.  If no failure witness exists, the search need not
terminate; this is a semi-decision procedure, not a decision procedure.

This semantic non-effectivity does not obstruct the identification-in-the-limit
result.  The learner receives neither the target grammar \(G\) nor a proof that
the target is \((f,h)\)-tuple-substitutable.  The theorem instead states that,
whenever the unknown target language actually belongs to \(\Cmcf{f}{h}\), there
exists a finite characteristic sample for that target, and after the learner has
seen any positive sample containing it, the hypothesis language equals the
target language.

If the target language is not \((f,h)\)-tuple-substitutable, the exact
reconstruction and identification theorems do not apply.  Thus the language \(L(\widehat G_0(K))\) constructed from a finite positive
sample \(K\) is not guaranteed to equal the
unknown target language.  Nevertheless, the learner is defined on every finite
positive sample, independently of whether the semantic promise holds, and its
hypothesis always generates every observed example,
\(K\subseteq L(\widehat G_0(K))\).  Hence the learner never discards an observed positive example.  The possible
failure concerns unobserved strings: the set
\(L(\widehat G_0(K))\setminus K\) may contain strings that are not in the target language.  Thus
\((f,h)\)-tuple substitutability is not needed to reproduce the observed
sample; it is what guarantees the soundness of extrapolation from the finite
sample to previously unseen strings.
\end{remark}


\section{Examples}
\label{sec:running-examples}

Before the abstract reconstruction proof, we record concrete grammars showing
what the fixed observation morphism is meant to capture in familiar block
synchronization languages.  The fixed regular envelope supplies only the coarse
zone pattern; the equality constraints are recovered distributionally from
positive examples.  We also verify the semantic promise for these examples, so
that the examples really belong to the relevant fixed-observation fibers.  For
related letter-count separations and calibrations in MCFLs, see
\cite{LehnerLindorfer2020}.

For a finite ordered alphabet \(\Sigma_m=\{a_1,\ldots,a_m\}\), put
\[
  R_m=a_1^*a_2^*\cdots a_m^* .
\]
Let \(D_m\) be the standard deterministic automaton for \(R_m\) with states
\(1,\ldots,m,\bot\), initial state \(1\), accepting states \(1,\ldots,m\), and
transitions
\[
  i\cdot a_j=
  \begin{cases}
    j, & i\le j,\\
    \bot, & i>j,
  \end{cases}
  \qquad
  \bot\cdot a_j=\bot .
\]
Let \(h_m:\Sigma_m^*\to M_m\) be its transition morphism.

\begin{lemma}[Zone morphism]
\label{lem:zone-morphism-basic}
For all words \(u,v\in\Sigma_m^*\), if \(h_m(u)=h_m(v)\), then
\(u\in R_m\) iff \(v\in R_m\).  Moreover, if \(m\ge2\) and
\(u\in \{a_1^n\cdots a_m^n\mid n\ge1\}\), then
\(h_m(u)\ne h_m(\eps)\).
\end{lemma}

\begin{proof}
The transition value \(h_m(w)\) determines the image of the initial state and
therefore determines membership in the regular language \(R_m\).  If \(m\ge2\) and
\(u=a_1^n\cdots a_m^n\) with \(n\ge1\), then the transition induced by \(u\)
sends the initial state \(1\) to \(m\), whereas the empty word induces the
identity transition and sends \(1\) to \(1\).  Hence the two transition values
are distinct.
\end{proof}

\begin{proposition}[Equality blocks are fixed-zone substitutable]
\label{prop:equality-block-substitutable}
For every \(m\ge2\), the language
\[
  E_m^+=\{a_1^n\cdots a_m^n\mid n\ge1\}
\]
is \((r,h_m)\)-tuple-substitutable for every arity \(r\ge1\).
\end{proposition}

\begin{proof}
Fix an arity \(d\ge1\).  Let
\(\tuple x,\tuple y\in(\Sigma_m^*)^d\) have the same componentwise
\(h_m\)-type, and suppose that they share an accepting sentence context \(E\):
\[
  E[\tuple x]=a_1^n\cdots a_m^n,
  \qquad
  E[\tuple y]=a_1^{n'}\cdots a_m^{n'}
\]
for some \(n,n'\ge1\).

For each \(1\le j\le m\), let \(k_j\) be the number of occurrences of \(a_j\)
in the fixed terminal part of \(E\), outside its named holes.  Since the same
sentence context \(E\) is used in both fillings, the value \(k_j\) is
independent of whether \(\tuple x\) or \(\tuple y\) is inserted.  Counting the
occurrences of \(a_j\) in the two displayed words gives
\[
  k_j+\sum_{i=1}^{d}|x_i|_{a_j}=n
\]
and
\[
  k_j+\sum_{i=1}^{d}|y_i|_{a_j}=n'.
\]
Subtracting the first equality from the second cancels the contribution of the
outside context and yields
\[
  \sum_{i=1}^{d}|y_i|_{a_j}
  -
  \sum_{i=1}^{d}|x_i|_{a_j}
  =n'-n
  \qquad (1\le j\le m).
  \tag{1}
\]

Now let \(F\) be any accepting sentence context for \(\tuple x\), say
\[
  F[\tuple x]=a_1^N\cdots a_m^N
\]
with \(N\ge1\).  By hypothesis,
\[
  h_m(x_i)=h_m(y_i)
  \qquad (1\le i\le d).
\]
The fixed terminal parts of \(F\) are unchanged, and the corresponding tuple
components are inserted into the same named holes.  Since \(h_m\) is a monoid
morphism, the value of a filled sentence context is the ordered product of the
values of its fixed terminal factors and inserted components.  Hence
\[
  h_m\bigl(F[\tuple x]\bigr)
  =
  h_m\bigl(F[\tuple y]\bigr).
\]
Because \(F[\tuple x]\in R_m\), Lemma~\ref{lem:zone-morphism-basic} implies
\[
  F[\tuple y]\in R_m.
\]

Equation~(1) shows that the number of occurrences of every letter \(a_j\) in
\(F[\tuple y]\) is \(N+(n'-n)\).  This common value cannot be zero.  Otherwise
\(F[\tuple y]=\eps\), and therefore
\[
  h_m(F[\tuple x])=h_m(F[\tuple y])=h_m(\eps),
\]
contradicting Lemma~\ref{lem:zone-morphism-basic}.  Thus
\(F[\tuple y]\in E_m^+\).  The reverse inclusion follows by exchanging
\(\tuple x\) and \(\tuple y\), so their tuple distributions are equal.
\end{proof}

The preceding argument can be generalized.  In particular, the framework below
covers languages such as
\[
  \{a^n\#_2b^m\#_3c^n\#_4d^m\mid n,m\ge1\}.
\]

\begin{definition}[Fixed ordered block envelope]
\label{def:fixed-ordered-block-envelope}
Let \(k\ge1\).  Take pairwise distinct letters
\[
  A=\{a_1,\ldots,a_k\}
\]
and a finite alphabet \(\Delta\) disjoint from \(A\).  Fix separator words
\[
  s_0,s_1,\ldots,s_k\in\Delta^*
\]
and put
\[
  \Sigma=A\uplus\Delta.
\]
The language
\[
  R
  =
  s_0a_1^+s_1a_2^+s_2\cdots s_{k-1}a_k^+s_k
  \subseteq\Sigma^*
\]
is called a \emph{fixed ordered block envelope}.

Every \(w\in R\) has a unique representation
\[
  w
  =
  s_0a_1^{n_1}s_1a_2^{n_2}s_2
  \cdots
  s_{k-1}a_k^{n_k}s_k,
  \qquad
  n_1,\ldots,n_k\ge1.
\]
Define its block-length vector by
\[
  \ell_R(w)
  :=
  (n_1,\ldots,n_k)
  \in\mathbb N_{>0}^k,
\]
and write \(\ell_j(w):=n_j\) for its \(j\)-th component.
\end{definition}

\begin{lemma}[Preservation of length equalities in fixed ordered block envelopes]
\label{lem:block-equality-envelopes}
Let \(R\subseteq\Sigma^*\) be a fixed ordered block envelope as in
Definition~\ref{def:fixed-ordered-block-envelope}.  Let \(\mathcal A_R\) be a
complete DFA recognizing \(R\), and let
\[
  h_R:\Sigma^*\longrightarrow M_R
\]
be its transition-monoid morphism.

For a finite set of pairs of block indices
\[
  \Lambda\subseteq\{1,\ldots,k\}^2,
\]
define
\[
  L_\Lambda
  :=
  \left\{
    w\in R
    \;\middle|\;
    \forall (p,q)\in\Lambda,\ 
    \ell_p(w)=\ell_q(w)
  \right\}.
\]
Then \(L_\Lambda\) is \((f,h_R)\)-tuple-substitutable for every \(f\ge1\).
\end{lemma}

\begin{proof}
Fix \(1\le d\le f\).  Let
\[
  \tuple x=(x_1,\ldots,x_d),
  \qquad
  \tuple y=(y_1,\ldots,y_d)
  \in(\Sigma^*)^d
\]
satisfy
\[
  h_R(x_i)=h_R(y_i)
  \qquad (1\le i\le d).
  \tag{2}
\]
Suppose further that \(\tuple x\) and \(\tuple y\) share an accepting
arity-\(d\) sentence context for \(L_\Lambda\).  Thus there is an arity-\(d\)
sentence context \(E\) such that
\[
  E[\tuple x]\in L_\Lambda,
  \qquad
  E[\tuple y]\in L_\Lambda.
  \tag{3}
\]

For \(1\le j\le k\) and
\(\tuple z=(z_1,\ldots,z_d)\), put
\[
  N_j(\tuple z)
  :=
  \sum_{i=1}^d |z_i|_{a_j}.
\]
This is the total number of occurrences of the visible block letter \(a_j\)
contributed by the components of \(\tuple z\).  Let \(C_j(E)\) be the number of
occurrences of \(a_j\) in the fixed terminal part of \(E\), outside its holes.
Then
\[
  |E[\tuple z]|_{a_j}
  =
  C_j(E)+N_j(\tuple z)
  \tag{4}
\]
for every \(\tuple z\).

Fix \((p,q)\in\Lambda\).  Since both fillings in~(3) belong to
\(L_\Lambda\), their \(p\)-th and \(q\)-th block lengths are equal.  Because
the separator alphabet is disjoint from the visible block letters, these
lengths are exactly the corresponding letter counts.  Hence
\[
  C_p(E)+N_p(\tuple x)
  =
  C_q(E)+N_q(\tuple x)
\]
and
\[
  C_p(E)+N_p(\tuple y)
  =
  C_q(E)+N_q(\tuple y).
\]
Subtracting gives
\[
  N_p(\tuple y)-N_p(\tuple x)
  =
  N_q(\tuple y)-N_q(\tuple x)
  \qquad ((p,q)\in\Lambda).
  \tag{5}
\]

Now take any
\[
  F\in\D_{L_\Lambda}^{(d)}(\tuple x),
\]
so that
\[
  F[\tuple x]\in L_\Lambda.
  \tag{6}
\]
By~(2) and the homomorphism property of \(h_R\), replacing each \(x_i\) by
\(y_i\) in the same named hole does not change the transition-monoid value of
the filled sentence context:
\[
  h_R(F[\tuple x])
  =
  h_R(F[\tuple y]).
  \tag{7}
\]
Since \(F[\tuple x]\in L_\Lambda\subseteq R\) and \(h_R\) is the transition
morphism of a DFA recognizing \(R\), equation~(7) implies
\[
  F[\tuple y]\in R.
  \tag{8}
\]

Let \(C_j(F)\) be the number of occurrences of \(a_j\) in the fixed terminal
part of \(F\).  For every \((p,q)\in\Lambda\), equation~(6) gives
\[
  C_p(F)+N_p(\tuple x)
  =
  C_q(F)+N_q(\tuple x).
  \tag{9}
\]
Combining~(5) and~(9), we obtain
\[
\begin{aligned}
  C_p(F)+N_p(\tuple y)
  &=
  C_p(F)+N_p(\tuple x)
    +\bigl(N_p(\tuple y)-N_p(\tuple x)\bigr)
  \\
  &=
  C_q(F)+N_q(\tuple x)
    +\bigl(N_q(\tuple y)-N_q(\tuple x)\bigr)
  \\
  &=
  C_q(F)+N_q(\tuple y).
\end{aligned}
\]
Therefore
\[
  \ell_p(F[\tuple y])
  =
  \ell_q(F[\tuple y])
  \qquad ((p,q)\in\Lambda).
\]
Together with~(8), this shows that
\[
  F[\tuple y]\in L_\Lambda.
\]
Hence
\[
  \D_{L_\Lambda}^{(d)}(\tuple x)
  \subseteq
  \D_{L_\Lambda}^{(d)}(\tuple y).
\]
The reverse inclusion follows by the same argument with \(\tuple x\) and
\(\tuple y\) exchanged.  Thus
\[
  \D_{L_\Lambda}^{(d)}(\tuple x)
  =
  \D_{L_\Lambda}^{(d)}(\tuple y).
\]
Since \(d\le f\) was arbitrary, \(L_\Lambda\) is
\((f,h_R)\)-tuple-substitutable.
\end{proof}

\begin{example}[The three-block agreement language]
\label{ex:l3-language}
Let
\[
  L_3=\{a^n b^n c^n\mid n\ge1\},
  \qquad R_3=a^*b^*c^* .
\]
Let \(h_{R_3}\) be the transition morphism of the standard DFA recognizing
\(R_3\).  The morphism records the coarse block zone of a component, but not the
equality of the three block lengths.
\end{example}

\begin{proposition}[A working presentation for \(L_3\)]
\label{prop:l3-running-example}
The language \(L_3\) has a reduced working binary linear nondeleting MCFG
presentation of fan-out two.
\end{proposition}

\begin{proof}
Let \(V=\{S,T,A,A_a,A_b,A_c\}\), let \(\mu(A)=2\), and let all other
nonterminals have fan-out \(1\).  Use the rules
\[
  S\to T,\qquad A_a\to(a),\qquad A_b\to(b),\qquad A_c\to(c),
\]
\[
  A\to(x^1,y^1)(A_a,A_b),\qquad
  A\to(a x^1,b x^2 y^1)(A,A_c),
\]
\[
  T\to(x^1x^2y^1)(A,A_c).
\]
A direct induction gives
\(L_A=\{(a^n,b^n c^{n-1})\mid n\ge1\}\), and the top rule yields exactly
\(a^n b^n c^n\).  The grammar is in the working form of
Definition~\ref{def:working-mcfg}.
\end{proof}

\begin{corollary}[The three-block example belongs to the fixed-observation class]
\label{cor:l3-in-fixed-observation-class}
For every \(f\ge2\), \(L_3\in\Cmcf{f}{h_{R_3}}\).
\end{corollary}

\begin{proof}
The presentation is Proposition~\ref{prop:l3-running-example}.  The semantic
promise follows from Proposition~\ref{prop:equality-block-substitutable}
with \((a_1,a_2,a_3)=(a,b,c)\).
\end{proof}

\begin{example}[The four-block parallel agreement language]
\label{ex:parallel-language}
Let
\[
  L_{\parallel}=\{a^n\#b^n\#c^n\#d^n\mid n\ge1\},
  \qquad P=a^+\#b^+\#c^+\#d^+ .
\]
Let \(h_P\) be the transition morphism of the standard DFA recognizing \(P\).
Again, the finite observation supplies the regular envelope; the equality of
block lengths is distributional information recovered from positive examples.
\end{example}

\begin{proposition}[A working presentation for \(L_{\parallel}\)]
\label{prop:parallel-running-example}
The language \(L_{\parallel}\) has a reduced working binary linear nondeleting
MCFG presentation of fan-out two.
\end{proposition}

\begin{proof}
Let \(H\) derive \((\#)\), let \(A_a,A_b,A_d\) derive \((a),(b),(d)\),
respectively, and let \(A\) have fan-out two.  Use
\[
  A\to(x^1\#y^1,\,c\#d)(A_a,A_b)
\]
and
\[
  A\to(a x^1 b,\,c x^2 y^1)(A,A_d).
\]
Then \(A\) derives exactly \((a^n\#b^n,c^n\#d^n)\) for \(n\ge1\).  A top rule
\[
  T\to(x^1y^1x^2)(A,H),
\]
together with \(S\to T\), yields exactly \(L_{\parallel}\).
\end{proof}

\begin{corollary}[The parallel example belongs to the fixed-observation class]
\label{cor:parallel-in-fixed-observation-class}
For every \(f\ge2\), \(L_{\parallel}\in\Cmcf{f}{h_P}\).
\end{corollary}

\begin{proof}
The presentation is Proposition~\ref{prop:parallel-running-example}.  View
\[
  P=a^+\#b^+\#c^+\#d^+
\]
as the fixed ordered block envelope with visible block letters
\((a_1,a_2,a_3,a_4)=(a,b,c,d)\), separator alphabet
\(\Delta=\{\#\}\), and separator words
\[
  (s_0,s_1,s_2,s_3,s_4)=(\eps,\#,\#,\#,\eps).
\]
The semantic promise follows from
Lemma~\ref{lem:block-equality-envelopes} with
\[
  \Lambda=\{(1,2),(2,3),(3,4)\}.
\]
\end{proof}

\begin{example}[The cross-serial two-spine language]
\label{ex:cross-serial-language}
Let
\[
  L_{\times}=\{a^n b^m c^n d^m\mid n,m\ge1\},
  \qquad Q=a^+b^+c^+d^+ .
\]
Let \(h_Q\) be the transition morphism of the standard DFA recognizing \(Q\).
This example has two independent agreement spines.
\end{example}

\begin{proposition}[A working presentation for \(L_{\times}\)]
\label{prop:cross-serial-running-example}
The language \(L_{\times}\) has a reduced working binary linear nondeleting MCFG
presentation of fan-out two.
\end{proposition}

\begin{proof}
Use two fan-out-two nonterminals \(A\) and \(B\).  Let terminal nonterminals
\(A_a,A_b,A_c,A_d\) derive \((a),(b),(c),(d)\), respectively.  The nonterminal
\(A\) derives \((a^n,c^n)\) by
\[
  A\to(x^1,y^1)(A_a,A_c),
  \qquad
  A\to(a x^1,x^2 y^1)(A,A_c),
\]
and \(B\) derives \((b^m,d^m)\) by
\[
  B\to(x^1,y^1)(A_b,A_d),
  \qquad
  B\to(b x^1,x^2 y^1)(B,A_d).
\]
The top rule
\[
  T\to(x^1y^1x^2y^2)(A,B),
\]
with \(S\to T\), generates exactly \(L_{\times}\).
\end{proof}

\begin{corollary}[The cross-serial example belongs to the fixed-observation class]
\label{cor:cross-serial-in-fixed-observation-class}
For every \(f\ge2\), \(L_{\times}\in\Cmcf{f}{h_Q}\).
\end{corollary}

\begin{proof}
The presentation is Proposition~\ref{prop:cross-serial-running-example}.  View
\[
  Q=a^+b^+c^+d^+
\]
as the fixed ordered block envelope with visible block letters
\((a_1,a_2,a_3,a_4)=(a,b,c,d)\), empty separator alphabet, and
\[
  s_0=s_1=s_2=s_3=s_4=\eps.
\]
The semantic promise follows from
Lemma~\ref{lem:block-equality-envelopes} with
\[
  \Lambda=\{(1,3),(2,4)\}.
\]
\end{proof}


\section{The Canonical Learner}
\label{sec:learner}

The learner uses only the fixed parameters \(f\) and \(h\) and the current finite
positive sample.  A witnessing target presentation is used only in the proof to
show that a finite characteristic sample exists.  The refinement below stores
only componentwise output \(h\)-types; the concrete exposing contexts needed for
reconstruction are supplied by the characteristic sample.

\subsection{Output-Type Refinement}
\label{sec:refinement}

The reconstruction proof uses a finite refinement of a witnessing grammar.  The
refinement records only the output \(h\)-type of each derived tuple.  It does not
record the order in which tuple components are placed in a surrounding sentence
or the \(h\)-values of boundary intervals.  These data are recovered from the
concrete witnessed segmentations and exposing contexts contained in the finite
characteristic sample.  This is sufficient because the learner's unit rule is
guarded by actual shared sample contexts.

\begin{definition}[Template evaluation]
\label{def:output-map}
Let
\[
  \rho\colon A\to(\alpha_1,\ldots,\alpha_e)(B,C)
\]
be a binary rule with \(\mu(B)=d_B\), \(\mu(C)=d_C\), and
\(e=\mu(A)\).  Put
\[
  \Gamma_\rho
  :=
  \Sigma
  \uplus
  \{x^1,\ldots,x^{d_B}\}
  \uplus
  \{y^1,\ldots,y^{d_C}\}.
\]
Each \(\alpha_\ell\) belongs to \(\Gamma_\rho^*\), and the elements of
\(\Gamma_\rho^*\) are called the \emph{template words} of the rule \(\rho\).

For
\[
  \mathbf q=(q_1,\ldots,q_{d_B})\in M^{d_B},
  \qquad
  \mathbf r=(r_1,\ldots,r_{d_C})\in M^{d_C},
\]
define
\[
  \eta_{\mathbf q,\mathbf r}\colon\Gamma_\rho\to M
\]
on letters by
\[
  \eta_{\mathbf q,\mathbf r}(a):=h(a)
  \qquad (a\in\Sigma),
\]
\[
  \eta_{\mathbf q,\mathbf r}(x^i):=q_i
  \qquad (1\le i\le d_B),
\]
and
\[
  \eta_{\mathbf q,\mathbf r}(y^j):=r_j
  \qquad (1\le j\le d_C).
\]
Let
\[
  \widehat{\eta}_{\mathbf q,\mathbf r}
  \colon
  \Gamma_\rho^*\to M
\]
be its unique extension to a monoid morphism.  Equivalently,
\[
  \widehat{\eta}_{\mathbf q,\mathbf r}(\eps)=1_M,
\]
and, for a template word
\[
  \beta=z_1\cdots z_t
  \qquad (z_1,\ldots,z_t\in\Gamma_\rho),
\]
\[
  \widehat{\eta}_{\mathbf q,\mathbf r}(\beta)
  =
  \eta_{\mathbf q,\mathbf r}(z_1)
  \cdots
  \eta_{\mathbf q,\mathbf r}(z_t).
\]

Define the evaluation of a template word \(\beta\) in the rule \(\rho\) by
\[
  \ev_\rho(\beta;\mathbf q,\mathbf r)
  :=
  \widehat{\eta}_{\mathbf q,\mathbf r}(\beta)
  \in M.
\]
Finally, define
\[
  \out_\rho(\mathbf q,\mathbf r)
  :=
  \bigl(
    \ev_\rho(\alpha_1;\mathbf q,\mathbf r),
    \ldots,
    \ev_\rho(\alpha_e;\mathbf q,\mathbf r)
  \bigr)
  \in M^e.
\]
\end{definition}

\begin{definition}[Output-type refinement]
\label{def:output-refinement}
Given \(G\) and \(h\), the full output-type refinement \(G^h\) has a fresh
start symbol \(\widetilde S\) and nonterminals \(A_{\mathbf p}\), where
\(A\in V\setminus\{S\}\), \(\mathbf p\in M^{\mu(A)}\).  Its rules are as follows.
For each start rule \(S\to A\) and each \(p\in M\), add
\(\widetilde S\to A_{(p)}\).  For each terminal rule \(A\to(a)\), add
\(A_{(h(a))}\to(a)\).  For each binary rule
\(\rho\colon A\to(\alpha_1,\ldots,\alpha_e)(B,C)\) and all
\(\mathbf q\in M^{\mu(B)}\), \(\mathbf r\in M^{\mu(C)}\), add
\(A_{\out_\rho(\mathbf q,\mathbf r)}\to\rho(B_{\mathbf q},C_{\mathbf r})\).
The trimmed output-type refinement \(\widetilde G_0\) is the subgrammar
obtained by keeping only typed nonterminals and typed rules that occur in some
successful derivation from \(\widetilde S\).
\end{definition}

\begin{proposition}[Output-type invariants]
\label{prop:output-invariants}
For the complete output-type refinement \(G^h\) and its trimmed subgrammar
\(\widetilde G_0\), the following statements hold.
\begin{enumerate}[label=(\roman*),leftmargin=*]
\item For every \(A\in V\setminus\{S\}\),
      \(\mathbf p\in M^{\mu(A)}\), and
      \(\tuple u\in(\Sigma^*)^{\mu(A)}\), if
      \(A_{\mathbf p}\Rightarrow_{G^h}^*\tuple u\), then
      \(h^{(\mu(A))}(\tuple u)=\mathbf p\).  Moreover, erasing all type
      indices from this \(G^h\)-derivation tree yields a valid
      \(G\)-derivation tree witnessing \(A\Rightarrow_G^*\tuple u\).

\item Every fixed \(G\)-derivation tree rooted at a nonterminal
      \(A\in V\setminus\{S\}\) lifts uniquely to a \(G^h\)-derivation tree by
      labeling each node with the componentwise \(h\)-type of the tuple derived
      at that node.  Here uniqueness means that, once the underlying
      \(G\)-derivation tree, including its rule labels and tree structure, is
      fixed, the type index at every node is uniquely determined.

\item \(L(G^h)=L(G)\), and, moreover,
      \(L(\widetilde G_0)=L(G^h)=L(G)\).
\end{enumerate}
\end{proposition}

\begin{proof}
We first relate template evaluation to actual string substitution in a binary
rule.  Let
\(\rho\colon A\to(\alpha_1,\ldots,\alpha_e)(B,C)\) be a binary rule, put
\(d_B=\mu(B)\) and \(d_C=\mu(C)\), and let
\(\tuple u=(u_1,\ldots,u_{d_B})\in(\Sigma^*)^{d_B}\) and
\(\tuple v=(v_1,\ldots,v_{d_C})\in(\Sigma^*)^{d_C}\).  Set
\(\mathbf q:=h^{(d_B)}(\tuple u)\) and
\(\mathbf r:=h^{(d_C)}(\tuple v)\).
For every template word \(\beta\in\Gamma_\rho^*\),
\[
  h\bigl(
    \widehat{\sigma}_{\tuple u,\tuple v}(\beta)
  \bigr)
  =
  \widehat{\eta}_{\mathbf q,\mathbf r}(\beta)
  =
  \ev_\rho(\beta;\mathbf q,\mathbf r).
  \tag{1}
\]
Indeed, the maps
\(h\circ\widehat{\sigma}_{\tuple u,\tuple v}\) and
\(\widehat{\eta}_{\mathbf q,\mathbf r}\colon\Gamma_\rho^*\to M\) are both
monoid morphisms.  They agree on every generator: for \(a\in\Sigma\),
\(h(\widehat{\sigma}_{\tuple u,\tuple v}(a))=h(a)
=\eta_{\mathbf q,\mathbf r}(a)\); for \(1\le i\le d_B\),
\(h(\widehat{\sigma}_{\tuple u,\tuple v}(x^i))=h(u_i)=q_i
=\eta_{\mathbf q,\mathbf r}(x^i)\); and, for \(1\le j\le d_C\),
\(h(\widehat{\sigma}_{\tuple u,\tuple v}(y^j))=h(v_j)=r_j
=\eta_{\mathbf q,\mathbf r}(y^j)\).
Equation~\textup{(1)} therefore follows from uniqueness of the morphism out of
the free monoid \(\Gamma_\rho^*\).

\medskip
\noindent
\textbf{Proof of (i).}
We argue by induction on the height of a \(G^h\)-derivation tree rooted at
\(A_{\mathbf p}\).

\smallskip
\noindent
\emph{Base case.}
If the tree has height one, its root rule is the terminal rule
\(A_{(h(a))}\to(a)\).  Hence \(\tuple u=(a)\),
\(\mathbf p=(h(a))\), and, since \(\mu(A)=1\),
\(h^{(\mu(A))}(\tuple u)=(h(a))=\mathbf p\).
Erasing the type index yields the original terminal rule \(A\to(a)\).

\smallskip
\noindent
\emph{Induction step.}
Suppose that the root rule is
\[
  A_{\mathbf p}
  \to
  \rho(B_{\mathbf q},C_{\mathbf r}).
  \tag{2}
\]
The underlying rule of \(G\) is
\(\rho\colon A\to(\alpha_1,\ldots,\alpha_e)(B,C)\), and, by
Definition~\ref{def:output-refinement},
\[
  \mathbf p=\out_\rho(\mathbf q,\mathbf r).
  \tag{3}
\]
Let the two child subtrees derive
\(B_{\mathbf q}\Rightarrow_{G^h}^*\tuple u\) and
\(C_{\mathbf r}\Rightarrow_{G^h}^*\tuple v\).
By the induction hypothesis,
\[
  h^{(\mu(B))}(\tuple u)=\mathbf q,
  \qquad
  h^{(\mu(C))}(\tuple v)=\mathbf r.
  \tag{4}
\]
Moreover, erasing type indices from the child subtrees yields valid
\(G\)-derivations \(B\Rightarrow_G^*\tuple u\) and
\(C\Rightarrow_G^*\tuple v\).
Let
\[
  \tuple w
  :=
  \rho(\tuple u,\tuple v)
  =
  \bigl(
    \widehat{\sigma}_{\tuple u,\tuple v}(\alpha_1),
    \ldots,
    \widehat{\sigma}_{\tuple u,\tuple v}(\alpha_e)
  \bigr)
\]
be the tuple derived at the root.  By equations~\textup{(1)}, \textup{(3)},
and~\textup{(4)}, for every \(1\le\ell\le e\),
\[
\begin{aligned}
  h(w_\ell)
  &=
  h\bigl(
    \widehat{\sigma}_{\tuple u,\tuple v}(\alpha_\ell)
  \bigr)
  \\
  &=
  \ev_\rho(\alpha_\ell;\mathbf q,\mathbf r).
\end{aligned}
\]
Consequently,
\[
\begin{aligned}
  h^{(e)}(\tuple w)
  &=
  \bigl(
    \ev_\rho(\alpha_1;\mathbf q,\mathbf r),
    \ldots,
    \ev_\rho(\alpha_e;\mathbf q,\mathbf r)
  \bigr)
  \\
  &=
  \out_\rho(\mathbf q,\mathbf r)
  \\
  &=
  \mathbf p.
\end{aligned}
\]
Erasing type indices from rule~\textup{(2)} gives the underlying rule \(\rho\)
of \(G\).  Placing this rule above the two erased child derivations yields a
valid \(G\)-derivation \(A\Rightarrow_G^*\tuple w\).
This proves~(i).

\medskip
\noindent
\textbf{Proof of (ii).}
Fix a \(G\)-derivation tree \(T\) rooted at
\(A\in V\setminus\{S\}\).  For each node \(\nu\), let \(A_\nu\) be its
nonterminal label and let
\(\tuple w_\nu\in(\Sigma^*)^{\mu(A_\nu)}\) be the tuple derived by the
subtree rooted at \(\nu\).  Define
\(\mathbf p_\nu:=h^{(\mu(A_\nu))}(\tuple w_\nu)\), and replace the node label
\(A_\nu\) by \((A_\nu)_{\mathbf p_\nu}\).
We show by induction on subtree height that this labeling produces a valid
\(G^h\)-derivation tree.

If \(\nu\) is a leaf using the terminal rule \(A_\nu\to(a)\), then
\(\tuple w_\nu=(a)\) and \(\mathbf p_\nu=(h(a))\), so \(G^h\) contains the
rule \((A_\nu)_{\mathbf p_\nu}\to(a)\).

Now suppose that an internal node \(\nu\) uses
\(\rho\colon A_\nu\to(\alpha_1,\ldots,\alpha_e)(B,C)\).
Let \(\tuple u\) and \(\tuple v\) be the tuples derived by its two child
subtrees, and put \(\mathbf q:=h^{(\mu(B))}(\tuple u)\) and
\(\mathbf r:=h^{(\mu(C))}(\tuple v)\).  Equation~\textup{(1)} gives
\(\mathbf p_\nu=\out_\rho(\mathbf q,\mathbf r)\).  Hence
Definition~\ref{def:output-refinement} provides the typed rule
\((A_\nu)_{\mathbf p_\nu}\to\rho(B_{\mathbf q},C_{\mathbf r})\).
Together with the induction hypotheses for the child subtrees, this proves
that the labeled tree is a valid \(G^h\)-derivation tree.

For uniqueness, fix the tree structure of \(T\), the rule used at each node,
and hence the tuple \(\tuple w_\nu\) derived at every node.  In any
\(G^h\)-lifting of \(T\), part~(i) forces the type index at \(\nu\) to be
\(h^{(\mu(A_\nu))}(\tuple w_\nu)=\mathbf p_\nu\).
Thus every node index, and therefore the entire lifting of \(T\), is unique.

\medskip
\noindent
\textbf{Proof of (iii).}
We first prove \(L(G^h)\subseteq L(G)\).
Let \(w\in L(G^h)\).  Then there are
\(A\in V\setminus\{S\}\) and \(p\in M\) such that a successful derivation
uses the root rule \(\widetilde S\to A_{(p)}\) and has a child subtree
witnessing \(A_{(p)}\Rightarrow_{G^h}^*(w)\).  By part~(i), erasing the type
indices from this subtree yields \(A\Rightarrow_G^*(w)\).  The typed start
rule was introduced from an original start rule \(S\to A\), so
\(S\Rightarrow_G^*(w)\).
Hence \(w\in L(G)\).

Conversely, let \(w\in L(G)\).  By start separation, a successful
\(G\)-derivation of \(w\) has a root start rule \(S\to A\) and, below it, a
derivation tree \(A\Rightarrow_G^*(w)\) whose root is not the start
symbol.  By part~(ii), this subtree lifts uniquely to
\(A_{(h(w))}\Rightarrow_{G^h}^*(w)\).  The grammar \(G^h\) contains the typed
start rule \(\widetilde S\to A_{(h(w))}\), so \(w\in L(G^h)\).  Therefore
\(L(G^h)=L(G)\).

Finally, \(\widetilde G_0\) is a subgrammar of \(G^h\), and hence
\(L(\widetilde G_0)\subseteq L(G^h)\).
For the reverse inclusion, let \(w\in L(G^h)\) and fix a successful
\(G^h\)-derivation tree \(T\) for \(w\).  By definition,
\(\widetilde G_0\) retains every typed nonterminal and typed rule that appears
in some successful derivation.  Every nonterminal and rule occurring in \(T\)
appears in the successful derivation \(T\) itself and is therefore retained.
Thus the same tree \(T\) is a successful \(\widetilde G_0\)-derivation, so
\(w\in L(\widetilde G_0)\).  It follows that
\(L(G^h)\subseteq L(\widetilde G_0)\).  Therefore
\(L(\widetilde G_0)=L(G^h)=L(G)\).
\end{proof}

For a surviving typed nonterminal \(X=A_{\mathbf p}\) in
\(\widetilde G_0\), put \(d:=\mu(A)\) and define its tuple language by
\[
  L_X
  :=
  \left\{
    \tuple u\in(\Sigma^*)^d
    \;\middle|\;
    X\Rightarrow_{\widetilde G_0}^*\tuple u
  \right\}.
\]
Fix once and for all a total effective order on tuples and on concrete sentence
contexts that refines length-lexicographic order.  Since \(X\) survives the
trimming, \(L_X\ne\emptyset\); let \(\omega(X)\) be the least tuple in \(L_X\).

\begin{lemma}[Concrete exposing contexts]
\label{lem:exposing-contexts}
Let \(X=A_{\mathbf p}\) be a typed nonterminal surviving in the trimmed
output-type refinement \(\widetilde G_0\), and put \(d:=\mu(A)\).
Then there exists a concrete arity-\(d\) sentence context \(E_X\) such that
\(E_X[\tuple u]\in L(G)\) for every \(\tuple u\in L_X\).  Moreover, if
\(\widetilde S\to X\) is a typed start rule of \(\widetilde G_0\), then
\(d=1\) and one may choose \(E_X=\blank_1\).
\end{lemma}

\begin{proof}
Because \(X\) survives the trimming, there is a successful derivation tree of
\(\widetilde G_0\) containing an occurrence of \(X\).  Fix such a successful
derivation tree \(T\), and fix a node \(\nu\) of \(T\) labeled by \(X\).  Let
\(T_\nu\) be the subtree rooted at \(\nu\).  This subtree derives some tuple
\(\tuple w=(w_1,\ldots,w_d)\in L_X\).
We keep fixed the portion of the derivation above \(\nu\) and all sibling
subtrees along the path from \(\nu\) to the root, and replace only the output of
\(T_\nu\) by formal named holes.

\medskip
\noindent
\textbf{Construction of the outside sentence context.}
Take fresh symbols \(\blank_1,\ldots,\blank_d\) not belonging to \(\Sigma\).
Delete the subtree \(T_\nu\) and place the formal tuple
\((\blank_1,\ldots,\blank_d)\) at node \(\nu\).  At every ancestor of \(\nu\), retain the rule used in \(T\),
and for every sibling subtree retain the concrete terminal tuple derived by
that subtree in \(T\).  Propagating the formal tuple from \(\nu\) to the root
by the simultaneous-substitution semantics of the intervening binary rules
produces, because the start symbol has fan-out one, a single word
\[
  E_{T,\nu}
  \in
  \bigl(
    \Sigma\cup\{\blank_1,\ldots,\blank_d\}
  \bigr)^*.
\]

Each hole \(\blank_i\) occurs exactly once in \(E_{T,\nu}\).  At node \(\nu\),
each hole occurs exactly once in the formal tuple
\((\blank_1,\ldots,\blank_d)\).  Every binary rule of the working MCFG is
linear and nondeleting: each variable corresponding to a child component
occurs exactly once in the entire parent template tuple.  Therefore, if every
hole occurs exactly once in the tuple at one level, then after substitution
into the parent rule every hole still occurs exactly once in the resulting
parent tuple.  Induction along the path from \(\nu\) to the root proves the
claim.

Consequently, there are a permutation \(\pi\) of \(\{1,\ldots,d\}\) and
words \(z_0,\ldots,z_d\in\Sigma^*\) such that
\[
  E_{T,\nu}
  =
  z_0\blank_{\pi(1)}z_1
  \cdots
  \blank_{\pi(d)}z_d.
\]
Thus \(E_{T,\nu}\) is an arity-\(d\) sentence context in the sense of
Definition~\ref{def:sentence-context}.  Set \(E_X:=E_{T,\nu}\).

\medskip
\noindent
\textbf{Substitution of an arbitrary \(X\)-derivation.}
Let \(\tuple u=(u_1,\ldots,u_d)\in L_X\).  By definition of \(L_X\), there
is a \(\widetilde G_0\)-derivation tree \(U\) witnessing
\(X\Rightarrow_{\widetilde G_0}^*\tuple u\).
Replace the subtree \(T_\nu\) of \(T\) by \(U\).  Since both subtrees have the
same root nonterminal \(X\), the resulting tree
\(T[\nu\leftarrow U]\) is again a valid successful \(\widetilde G_0\)-derivation tree.

The holes of \(E_X\) record where the components derived at \(\nu\) occur in
the final yield.  Hence
\[
  \yield\bigl(T[\nu\leftarrow U]\bigr)
  =
  E_X[\tuple u].
\]
Formally, this equality follows by induction on the length of the path from
\(\nu\) to the root: at each ancestor, the tuple obtained after replacing
\(\blank_i\) by \(u_i\) is exactly the tuple obtained by applying the same
binary rule to the replaced child derivation and the fixed sibling derivation.
At the root this gives the displayed equality.

Therefore \(E_X[\tuple u]\in L(\widetilde G_0)\).  By
Proposition~\ref{prop:output-invariants}, \(L(\widetilde G_0)=L(G)\), and
hence \(E_X[\tuple u]\in L(G)\).
Since \(\tuple u\in L_X\) was arbitrary, the required property holds for every
\(\tuple u\in L_X\).

Finally, suppose that \(\widetilde S\to X\) is a typed start rule of
\(\widetilde G_0\).  Such a rule comes from an
original start rule \(S\to A\), and the working-form definition gives
\(\mu(A)=1\).  Thus \(d=1\).  Choose a successful derivation tree whose root rule is
\(\widetilde S\to X\), and delete its unique child subtree rooted at \(X\).
No terminal material remains outside that subtree, and its sole component is
the entire generated string.  Therefore one may choose \(E_X=\blank_1\).
\end{proof}

\begin{definition}[Exposing context]
\label{def:exposing-context}
For each surviving \(X\), choose one successful derivation of
\(\widetilde G_0\) in which \(X\) occurs, and let \(\chi(X)\) be the concrete
sentence context obtained by deleting the subtree rooted at that occurrence and
keeping the named holes of the exposed tuple.  By
Lemma~\ref{lem:exposing-contexts}, this context satisfies
\(\chi(X)[\tuple u]\in L(G)\) for every \(\tuple u\in L_X\).  For a start child
\(X\), the convention chooses \(\chi(X)=\blank_1\).
\end{definition}

\subsection{The characteristic sample}
\label{subsec:learner-grammar}

\begin{definition}[Presentation-relative characteristic sample of a witnessing refinement]
\label{def:cs}
Let \(\widetilde G_0\) be the trimmed output-type refinement.  For each
surviving typed nonterminal \(X\), include the word \(\chi(X)[\omega(X)]\).  For
each terminal rule \(X\to(a)\) of \(\widetilde G_0\), include
\(\chi(X)[(a)]\).  For each binary rule \(R\colon X\to\rho(Y,Z)\) of
\(\widetilde G_0\), include \(\chi(X)[\rho(\omega(Y),\omega(Z))]\).  The union
of these finitely many words is denoted \(\CS(\widetilde G_0)\).
\end{definition}

The presentation-relative sample \(\CS(\widetilde G_0)\) is relative to the chosen
witnessing presentation \(G\) and its trimmed output-type refinement.  The learner is not
given this presentation or the sample.  The role of \(\CS(\widetilde G_0)\) is
only to prove the existence of a finite sufficient positive set in the sense of
Gold identification.

\begin{lemma}[The characteristic sample is positive and finite]
\label{lem:cs-positive}
\(\CS(\widetilde G_0)\) is a finite subset of \(L(G)\).
\end{lemma}

\begin{proof}
There are finitely many surviving typed nonterminals and typed rules.  By
Definition~\ref{def:exposing-context}, \(\chi(X)\) accepts every tuple in
\(L_X\).  The anchor \(\omega(X)\), the terminal tuple \((a)\) for a terminal
rule \(X\to(a)\), and the tuple \(\rho(\omega(Y),\omega(Z))\) for a binary rule
\(X\to\rho(Y,Z)\) all belong to \(L_X\).  Thus each displayed word lies in
\(L(G)\).
\end{proof}

\begin{proposition}[Concrete-context sufficiency]
\label{prop:concrete-context-sufficiency}
Let \(G\) be a witnessing working presentation and let \(\widetilde G_0\) be
its trimmed output-type refinement.  For each surviving typed nonterminal
\(X\), choose an anchor tuple \(\omega(X)\) and a concrete exposing context
\(\chi(X)\) that accepts every tuple derived from \(X\).  Then the data
consisting of
\begin{enumerate}[label=(\roman*),leftmargin=*]
\item the componentwise output \(h\)-type of each surviving typed symbol,
\item one exposed anchor for each surviving typed symbol, and
\item one exposed filled instance of each surviving typed rule
\end{enumerate}
are sufficient for the canonical learner to simulate every rule of
\(\widetilde G_0\).  No additional sentence-level interface label is required
on the learner's nonterminals.
\end{proposition}

\begin{proof}
The output type guarantees that a typed parent anchor and the corresponding
filled rule tuple have the same componentwise \(h\)-value.  Their two exposed
sample words use the same concrete context \(\chi(X)\), so the learner may
connect the two observed tuple values by a guarded unit rule.  A filled binary
rule word records the actual variable intervals of its two child anchors inside
the parent tuple; the anchors themselves may be observed through their own
exposing words.  The resulting binary witness reconstructs the underlying
linear template.  Terminal and start rules are handled by the corresponding
exposed terminal and anchor words.  The formal rule-by-rule simulation is
Lemma~\ref{lem:rule-simulation}.
\end{proof}

\subsection{Observed tuples and concrete witnesses}

\begin{definition}[Tuple occurrence in a sample word]
\label{def:tuple-occurrence}
Let \(K\subseteq\Sigma^*\), let \(w=a_1a_2\cdots a_n\in K\), and let
\(d\ge1\).  We regard \(0,1,\ldots,n\) as the cut positions of \(w\).  For
\(0\le p\le q\le n\), put
\(w[p:q]:=a_{p+1}a_{p+2}\cdots a_q\), with \(w[p:p]:=\eps\).

An \emph{arity-\(d\) tuple occurrence} in \(w\) is data
\[
  \mathfrak o=\bigl(\sigma,(\ell_i,r_i)_{i=1}^d\bigr),
\]
where \(\sigma\) is a permutation of \(\{1,\ldots,d\}\), each
\(\ell_i,r_i\) is a cut position of \(w\), and
\[
  0\le \ell_{\sigma(1)}\le r_{\sigma(1)}\le
  \ell_{\sigma(2)}\le r_{\sigma(2)}\le\cdots\le
  \ell_{\sigma(d)}\le r_{\sigma(d)}\le n.
\]
The associated tuple is
\(\tuple x_{\mathfrak o}=(x_1,\ldots,x_d)\in(\Sigma^*)^d\), where
\(x_i:=w[\ell_i:r_i]\).

Define \(u_0,\ldots,u_d\in\Sigma^*\) by
\[
  u_0:=w[0:\ell_{\sigma(1)}],\qquad
  u_j:=w[r_{\sigma(j)}:\ell_{\sigma(j+1)}]\ (1\le j<d),\qquad
  u_d:=w[r_{\sigma(d)}:n].
\]
Then
\[
  E_{\mathfrak o}:=
  u_0\blank_{\sigma(1)}u_1\blank_{\sigma(2)}\cdots
  u_{d-1}\blank_{\sigma(d)}u_d
\]
is an arity-\(d\) sentence context and
\(E_{\mathfrak o}[\tuple x_{\mathfrak o}]=w\).  We therefore identify
\(\mathfrak o\) with the pair
\((E_{\mathfrak o},\tuple x_{\mathfrak o})\).  Equivalently, an arity-\(d\)
tuple occurrence in \(w\) is a pair \((E,\tuple x)\) such that \(E\) is an
arity-\(d\) sentence context, \(\tuple x\in(\Sigma^*)^d\), and
\(E[\tuple x]=w\).

For each \(i\), if \(\ell_i<r_i\), the \(i\)-th slot is the positive-length
half-open interval \([\ell_i,r_i)\); if \(\ell_i=r_i\), it is an empty slot at
the cut position \(\ell_i\).  If distinct empty slots \(i,j\) share a cut
position \(c\), that is,
\(\ell_i=r_i=\ell_j=r_j=c\), then their local order is induced by \(\sigma\):
slot \(i\) precedes slot \(j\) exactly when
\(\sigma^{-1}(i)<\sigma^{-1}(j)\).
\end{definition}

\begin{example}[Empty slots and local tie order]
\label{ex:empty-slots}
Let \(w=ab\) and consider the tuple \(\tuple x=(a,\eps,b)\).  The context
\(E=\blank_1\blank_2\blank_3\) satisfies \(E[\tuple x]=ab\).  As an occurrence
inside \(w\), this is represented by the interval \([0,1]\) for
\(\blank_1\), the zero-length interval \([1,1]\) for \(\blank_2\), and the
interval \([1,2]\) for \(\blank_3\).  If two empty slots occur at the same cut,
the local tie order distinguishes them.  For instance, with
\(\tuple y=(a,\eps,\eps,b)\), both
\(\blank_1\blank_2\blank_3\blank_4\) and
\(\blank_1\blank_3\blank_2\blank_4\) fill to the same word \(ab\), but they are
different concrete occurrences.  The tie order at the cut between \(a\) and
\(b\) records whether \(\blank_2<\blank_3\) or \(\blank_3<\blank_2\).  This is
only bookkeeping, but it lets binary witnesses reconstruct a unique linear
template even when empty components occur.
\end{example}

\begin{definition}[Observed tuple]
\label{def:observed-tuple}
A tuple \(\tuple x\in(\Sigma^*)^d\) with \(1\le d\le f\) is \emph{observed in
\(K\)} if \((E,\tuple x)\) is a tuple occurrence in some word of \(K\).  The
learner uses one nonterminal \([\tuple x]\) for each observed tuple, together
with a fresh start symbol \(\widehat S\).
\end{definition}

\begin{definition}[Binary witness]
\label{def:binary-witness}
Let \(K\subseteq\Sigma^*\) be finite and let \(e,d_B,d_C\ge1\).  Put
\[
  X_B:=\{x^1,\ldots,x^{d_B}\},\qquad
  Y_C:=\{y^1,\ldots,y^{d_C}\},\qquad
  \Gamma:=\Sigma\uplus X_B\uplus Y_C.
\]
A \emph{binary witness} in \(K\), of output arity \(e\), left-child arity
\(d_B\), and right-child arity \(d_C\), is data
\[
  \mathfrak b=\bigl(
    \mathfrak o_P,\mathfrak o_B,\mathfrak o_C,(m_i)_{i=1}^e,
    (u_{i,r})_{\substack{1\le i\le e\\0\le r\le m_i}},
    (\xi_{i,r})_{\substack{1\le i\le e\\1\le r\le m_i}}
  \bigr)
\]
satisfying the following conditions.
\begin{enumerate}[label=(\roman*),leftmargin=*]
\item There is an arity-\(e\) tuple occurrence
      \(\mathfrak o_P=(E_P,\tuple z)\) in some \(w_P\in K\), where
      \(\tuple z=(z_1,\ldots,z_e)\in(\Sigma^*)^e\) and
      \(E_P[\tuple z]=w_P\).  It is called the \emph{parent occurrence}.

\item There are tuple occurrences
      \(\mathfrak o_B=(E_B,\tuple x)\) and
      \(\mathfrak o_C=(E_C,\tuple y)\) in words \(w_B,w_C\in K\), where
      \(\tuple x=(x_1,\ldots,x_{d_B})\in(\Sigma^*)^{d_B}\),
      \(\tuple y=(y_1,\ldots,y_{d_C})\in(\Sigma^*)^{d_C}\), and
      \(E_B[\tuple x]=w_B\), \(E_C[\tuple y]=w_C\).  The words
      \(w_P,w_B,w_C\) may be equal or distinct; in particular, the child
      observations need not occur in the parent sample word.

\item For each \(1\le i\le e\), an integer \(m_i\ge0\), terminal words
      \(u_{i,0},\ldots,u_{i,m_i}\in\Sigma^*\), and labels
      \(\xi_{i,1},\ldots,\xi_{i,m_i}\in X_B\uplus Y_C\) are specified.
      The map \((i,r)\mapsto\xi_{i,r}\) is a bijection from its index set
      onto \(X_B\uplus Y_C\).  Thus every variable occurs exactly once across
      the complete parent tuple.

\item Define \(\tau_{\tuple x,\tuple y}:X_B\uplus Y_C\to\Sigma^*\) by
      \(\tau_{\tuple x,\tuple y}(x^j):=x_j\) and
      \(\tau_{\tuple x,\tuple y}(y^k):=y_k\).  For every \(1\le i\le e\),
      \[
        z_i=u_{i,0}\tau_{\tuple x,\tuple y}(\xi_{i,1})u_{i,1}\cdots
        \tau_{\tuple x,\tuple y}(\xi_{i,m_i})u_{i,m_i}.
        \tag{*}
      \]
      Thus \textup{(*)} decomposes \(z_i\) into terminal gaps and variable
      intervals whose values are child-tuple components.

\item For \(1\le r\le m_i\), define the endpoints of the interval labelled
      \(\xi_{i,r}\) in \(z_i\) by
      \[
      \begin{aligned}
        \ell_{i,r}
        &:={}|u_{i,0}|+\sum_{q=1}^{r-1}
        \bigl(|\tau_{\tuple x,\tuple y}(\xi_{i,q})|+|u_{i,q}|\bigr),\\
        r_{i,r}
        &:={\ell_{i,r}}+|\tau_{\tuple x,\tuple y}(\xi_{i,r})|.
      \end{aligned}
      \]
      Then \(z_i[\ell_{i,r}:r_{i,r}]
      =\tau_{\tuple x,\tuple y}(\xi_{i,r})\), and
      \[
        0\le\ell_{i,1}\le r_{i,1}\le\cdots\le
        \ell_{i,m_i}\le r_{i,m_i}\le|z_i|.
      \]

\item If \(\tau_{\tuple x,\tuple y}(\xi_{i,r})=\eps\), the corresponding
      variable interval is the zero-length interval
      \([\ell_{i,r},\ell_{i,r})\).  If several such intervals occur at the
      same cut position \(c\), their local order is the order in the list
      \(\xi_{i,1},\ldots,\xi_{i,m_i}\): if
      \(\ell_{i,r}=r_{i,r}=\ell_{i,s}=r_{i,s}=c\) and \(r<s\), then the
      interval labelled \(\xi_{i,r}\) precedes that labelled \(\xi_{i,s}\).
\end{enumerate}

For \(1\le i\le e\), define
\(\alpha_i:=u_{i,0}\xi_{i,1}u_{i,1}\cdots
\xi_{i,m_i}u_{i,m_i}\in\Gamma^*\), and put
\(\boldsymbol\alpha_{\mathfrak b}:=(\alpha_1,\ldots,\alpha_e)\).
Condition~(iii) makes this a binary linear nondeleting template tuple.
Let \(\widehat\tau_{\tuple x,\tuple y}:\Gamma^*\to\Sigma^*\) be the unique
monoid-morphism extension of \(\tau_{\tuple x,\tuple y}\), acting identically
on \(\Sigma\).  Condition~(iv) says
\(\widehat\tau_{\tuple x,\tuple y}(\alpha_i)=z_i\) for every \(i\).  Hence the
associated template \(\rho_{\mathfrak b}\), defined by
\[
  \rho_{\mathfrak b}(\tuple x,\tuple y):=
  \bigl(\widehat\tau_{\tuple x,\tuple y}(\alpha_1),\ldots,
        \widehat\tau_{\tuple x,\tuple y}(\alpha_e)\bigr),
\]
satisfies \(\tuple z=\rho_{\mathfrak b}(\tuple x,\tuple y)\).
\end{definition}

\begin{remark}[Locality of a binary witness]
\label{rem:binary-witness-locality}
The parent occurrence supplies the positions at which the already observed
child tuple values are used.  It is not required to supply separate exposing
sentence contexts for the children inside the same sample word.  In the
completeness proof, the child anchors are observed through their own exposing
contexts, while the filled parent rule word supplies their concrete intervals
inside the parent tuple.
\end{remark}

\begin{lemma}[Uniqueness of the induced binary template]
\label{lem:unique-induced-template}
Let \(e,d_B,d_C\ge1\), put
\(X_B:=\{x^1,\ldots,x^{d_B}\}\),
\(Y_C:=\{y^1,\ldots,y^{d_C}\}\), and
\(\Gamma:=\Sigma\uplus X_B\uplus Y_C\).  Fix a parent tuple
\(\tuple z=(z_1,\ldots,z_e)\in(\Sigma^*)^e\) and child tuples
\(\tuple x=(x_1,\ldots,x_{d_B})\),
\(\tuple y=(y_1,\ldots,y_{d_C})\).

For every \(1\le i\le e\), suppose that an ordered segmentation is specified by
\[
  z_i=u_{i,0}\tau_{\tuple x,\tuple y}(\xi_{i,1})u_{i,1}\cdots
      \tau_{\tuple x,\tuple y}(\xi_{i,m_i})u_{i,m_i},
  \tag{1}
\]
where \(m_i\ge0\), \(u_{i,0},\ldots,u_{i,m_i}\in\Sigma^*\), and
\(\xi_{i,1},\ldots,\xi_{i,m_i}\in X_B\uplus Y_C\).  Here
\(\tau_{\tuple x,\tuple y}(x^j):=x_j\) and
\(\tau_{\tuple x,\tuple y}(y^k):=y_k\).  Assume that
\[
  \{\xi_{i,r}\mid 1\le i\le e,\ 1\le r\le m_i\}=X_B\uplus Y_C,
  \tag{2}
\]
and that \((i,r)\mapsto\xi_{i,r}\) is injective.  For zero-length intervals
sharing a cut position, also fix the local order given by the corresponding
label list.

Then there is a unique binary linear nondeleting template tuple
\(\boldsymbol\alpha=(\alpha_1,\ldots,\alpha_e)\in(\Gamma^*)^e\) consistent
with these ordered segmentations, namely
\[
  \alpha_i=u_{i,0}\xi_{i,1}u_{i,1}\cdots
  \xi_{i,m_i}u_{i,m_i}\qquad(1\le i\le e).
  \tag{3}
\]
If \(\widehat\tau_{\tuple x,\tuple y}:\Gamma^*\to\Sigma^*\) is the unique
monoid-morphism extension of \(\tau_{\tuple x,\tuple y}\), acting identically
on \(\Sigma\), then
\(\widehat\tau_{\tuple x,\tuple y}(\alpha_i)=z_i\) for every \(i\).
Consequently, the binary template \(\rho\) determined by
\(\boldsymbol\alpha\) satisfies
\(\rho(\tuple x,\tuple y)=\tuple z\).
\end{lemma}

\begin{proof}
Define \(\alpha_i\) by~\textup{(3)} and put
\(\boldsymbol\alpha:=(\alpha_1,\ldots,\alpha_e)\).  By~\textup{(2)} and the
injectivity of \((i,r)\mapsto\xi_{i,r}\), every variable occurs exactly once
in \(\boldsymbol\alpha\); hence the tuple is linear and nondeleting.  A
variable whose child component is empty still occurs as a symbol in the
template.  When several zero-length intervals share a cut position, their
specified local order fixes their order in the template word, so~\textup{(3)}
is unambiguous.

Because \(\widehat\tau_{\tuple x,\tuple y}\) fixes terminal letters and maps
each variable to its corresponding child component,~\textup{(1)} gives
\[
  \widehat\tau_{\tuple x,\tuple y}(\alpha_i)
  =u_{i,0}\tau_{\tuple x,\tuple y}(\xi_{i,1})u_{i,1}\cdots
   \tau_{\tuple x,\tuple y}(\xi_{i,m_i})u_{i,m_i}
  =z_i.
\]
Thus \(\rho(\tuple x,\tuple y)=\tuple z\).

For uniqueness, let
\(\boldsymbol\beta=(\beta_1,\ldots,\beta_e)\) be another binary linear
nondeleting template tuple consistent with the same ordered segmentations.
Consistency requires the terminal parts of \(\beta_i\), read left to right, to
be \(u_{i,0},\ldots,u_{i,m_i}\), and the intervening variables, in their
specified order, to be \(\xi_{i,1},\ldots,\xi_{i,m_i}\).  Therefore
\(\beta_i=u_{i,0}\xi_{i,1}u_{i,1}\cdots
\xi_{i,m_i}u_{i,m_i}=\alpha_i\) for every \(i\).  Hence
\(\boldsymbol\beta=\boldsymbol\alpha\).
\end{proof}

\begin{definition}[Canonical learner]
\label{def:learner}
The canonical learner grammar \(\widehat G_0(K)\) has start symbol
\(\widehat S\), nonterminals \([\tuple x]\) for observed tuples of arity at most
\(f\), and the following rules.
\begin{enumerate}[label=(\roman*),leftmargin=*]
\item For each \(w\in K\), add the start rule \(\widehat S\to[(w)]\).
\item For each observed unary tuple \((a)\) with \(a\in\Sigma\), add
\([(a)]\to(a)\).
\item For each binary witness with parent occurrence \((E,\tuple z)\), child
tuples \(\tuple x,\tuple y\), and induced template \(\rho\), add
\([\tuple z]\to\rho([\tuple x],[\tuple y])\).
\item For observed tuples \(\tuple x,\tuple y\) of the same arity \(d\), add the
unit rule \([\tuple x]\to[\tuple y]\) if \(h^{(d)}(\tuple x)=h^{(d)}(\tuple y)\)
and there exists a concrete arity-\(d\) context \(E\) such that
\(E[\tuple x]\in K\) and \(E[\tuple y]\in K\).
\end{enumerate}
The grammar is extended by unit rules.  They can be eliminated by the standard
unit-closure construction if a unit-free presentation is desired.
\end{definition}

\begin{remark}[Meaning of ``canonical'']
\label{rem:canonical-meaning}
Here ``canonical'' means deterministic and set-driven relative to fixed
effective orders on tuple values, occurrences, segmentations, and output
rules.  It does not mean that the returned grammar is a canonical minimal
presentation of the target language.
\end{remark}

\begin{remark}[Pseudocode for the canonical learner]
\label{rem:learner-pseudocode}
For clarity, the construction can be written as the following set-driven
procedure.
\begin{enumerate}[label=\textnormal{Step \arabic*.},leftmargin=*]
\item Enumerate all tuple occurrences \((E,\tuple x)\) of arity at most \(f\) in
the current finite sample \(K\), and create one nonterminal \([\tuple x]\) for
each observed tuple value.
\item Add start rules \(\widehat S\to[(w)]\) for all \(w\in K\), and add terminal
rules \([(a)]\to(a)\) for every observed unary letter tuple.
\item Enumerate all concrete binary witnesses inside sample words and add the
corresponding binary linear nondeleting rules.
\item For every pair of observed tuples \(\tuple x,\tuple y\) of the same arity,
add the unit rule \([\tuple x]\to[\tuple y]\) exactly when they have the same
componentwise \(h\)-type and are seen in one common concrete sentence context in
\(K\).
\item Output the resulting grammar, optionally after eliminating unit rules by
standard unit-closure.
\end{enumerate}
All loops range over explicitly enumerable cuts and intervals in the finite
sample.  Hence the procedure is effective, and its fixed-parameter polynomial
bound is proved in Theorem~\ref{thm:poly}.
\end{remark}

\begin{example}[A small trace of the fixed-\(h\) unit rule]
\label{ex:learner-trace}
Let \(P=a^+\#b^+\#c^+\#d^+\), let \(h_P\) be its transition morphism, and take
\[
  K=\{a\#b\#c\#d,\ aa\#bb\#cc\#dd\}.
\]
The tuple \(\tuple x=(a\#b,c\#d)\) occurs in the first sample word with context
\(E=\blank_1\#\blank_2\).  The tuple
\(\tuple y=(aa\#bb,cc\#dd)\) occurs in the second sample word with the same
context.  In the standard DFA for \(P\), the words \(a\#b\) and \(aa\#bb\)
induce the same transition, and the words \(c\#d\) and \(cc\#dd\) induce the
same transition.  Hence
\(h_P^{(2)}(\tuple x)=h_P^{(2)}(\tuple y)\), and the learner adds the unit rules
\([\tuple x]\to[\tuple y]\) and \([\tuple y]\to[\tuple x]\).  This is the
finite-type guarded version of the shared-context substitution step; without the
\(h_P\)-type check, analogous untyped shared-context triggers are too coarse
for such synchronization examples.
\end{example}

The learner receives only \(K\), \(f\), and \(h\).  It does not receive the
target grammar, the refinement, or the characteristic sample.


\begin{observation}[Set-driven consistency]
\label{obs:set-driven-consistent}
The learner $K\mapsto\widehat G_0(K)$ is set-driven.  Moreover, for the
$\eps$-free samples considered in this paper, it is consistent:
$K\subseteq L(\widehat G_0(K))$.
\end{observation}

\begin{proof}
Set-drivenness is immediate from Definition~\ref{def:learner}: all observed
tuples, witnesses, and unit rules are computed from the finite set $K$, not
from the order in which the words were presented.  For consistency, take
$w=a_1\cdots a_n\in K$ with $n\ge1$.  The learner has the start rule
$\widehat S\to[(w)]$.  Each letter $a_i$ occurs as an observed unary tuple, so
the terminal rule $[(a_i)]\to(a_i)$ is present.  For every nontrivial split
$uv$ of a substring of $w$, the parent tuple $(uv)$, the child tuples $(u)$ and
$(v)$, and the template $x^1y^1$ form a binary witness.  A binary induction on
the length of substrings therefore gives $[(w)]\Rightarrow^*(w)$.  Hence
$w\in L(\widehat G_0(K))$, and this holds for all $w\in K$.
\end{proof}

\begin{observation}[Monotonicity in the positive sample]
\label{obs:learner-monotone}
If \(K\subseteq K'\), then every nonterminal and every rule of
\(\widehat G_0(K)\) also occurs in \(\widehat G_0(K')\), up to the evident
inclusion of observed tuple values.  Hence
\(L(\widehat G_0(K))\subseteq L(\widehat G_0(K'))\).
\end{observation}

\begin{proof}
All four classes of rules in Definition~\ref{def:learner} are generated from
positive evidence present in the finite sample: sample words, observed unary
letter tuples, concrete binary witnesses, and shared-context evidence for unit
rules.  Passing from \(K\) to a larger sample \(K'\) can add such evidence but
cannot remove it.  The language inclusion follows because the grammar only gains
nonterminals and rules.
\end{proof}


\section{Soundness and Exact Reconstruction}
\label{sec:exact}

This section proves that the learner is both conservative with respect to the
semantic promise and complete once the sample contains the presentation-relative
characteristic sample.

\subsection{Soundness}
\label{sec:soundness}

\begin{definition}[Fixed-\(h\) distributional equivalence]
\label{def:distributional-equivalence}
For \(\tuple u,\tuple x\in(\Sigma^*)^d\), write
\(\tuple u\equiv_L^d\tuple x\) if
\(h^{(d)}(\tuple u)=h^{(d)}(\tuple x)\) and
\(\D_L^{(d)}(\tuple u)=\D_L^{(d)}(\tuple x)\).  The superscript is omitted when
clear.
\end{definition}

\begin{lemma}[Shared-context substitutability]
\label{lem:shared-context}
Let \(L\) be \((f,h)\)-tuple-substitutable.  If \(d\le f\),
\(h^{(d)}(\tuple x)=h^{(d)}(\tuple y)\), and there is a concrete arity-\(d\)
context \(E\) with \(E[\tuple x],E[\tuple y]\in L\), then
\(\tuple x\equiv_L^d\tuple y\).
\end{lemma}

\begin{proof}
The context \(E\) lies in both tuple distributions.  Definition~\ref{def:tuple-substitutable}
therefore gives equality of the distributions, and the \(h\)-type equality is
part of the hypothesis.
\end{proof}

\begin{lemma}[Filling identity]
\label{lem:filling-identity}
Let \(\rho\colon\bullet\to(\alpha_1,\ldots,\alpha_e)(\bullet,\bullet)\) be a
binary linear nondeleting template.  If \(E\) is an arity-\(e\) sentence context
around the parent tuple and \(\tuple y\) is fixed, then replacing the
\(x\)-variables in \(E[\rho(\tuple x,\tuple y)]\) by holes defines a concrete
arity-\(d_B\) context \(E_B^{\rho,E,\tuple y}\) satisfying
\(E_B^{\rho,E,\tuple y}[\tuple x]=E[\rho(\tuple x,\tuple y)]\).  The analogous
context \(E_C^{\rho,E,\tuple x}\) satisfies
\(E_C^{\rho,E,\tuple x}[\tuple y]=E[\rho(\tuple x,\tuple y)]\).
\end{lemma}

\begin{proof}
Each child variable occurs exactly once in the template.  Holding the sibling
tuple fixed and the parent context fixed, the positions formerly occupied by
the chosen child variables become named holes, with the inherited left-to-right
order.  Substituting the child tuple back into these holes recovers the original
filled sentence.
\end{proof}

\begin{lemma}[Witnessed composition preserves equivalence]
\label{lem:witnessed-composition}
Let \(L\) be \((f,h)\)-tuple-substitutable.  Suppose that
\(E[\rho(\tuple x,\tuple y)]\in L\) for a concrete parent context \(E\), and
that \(\tuple u\equiv_L\tuple x\) and \(\tuple v\equiv_L\tuple y\).  Then
\(\rho(\tuple u,\tuple v)\equiv_L\rho(\tuple x,\tuple y)\).
\end{lemma}

\begin{proof}
By Lemma~\ref{lem:filling-identity}, the induced \(B\)-side context satisfies
\(E_B^{\rho,E,\tuple y}[\tuple x]=E[\rho(\tuple x,\tuple y)]\in L\).  Since
\(\tuple u\equiv_L\tuple x\), this context also accepts \(\tuple u\), so
\(E[\rho(\tuple u,\tuple y)]\in L\).  Now use the induced \(C\)-side context
with sibling \(\tuple u\).  Since \(\tuple v\equiv_L\tuple y\), we get
\(E[\rho(\tuple u,\tuple v)]\in L\).  The componentwise \(h\)-types of the two
parent tuples are equal because \(h^{(d_B)}(\tuple u)=h^{(d_B)}(\tuple x)\),
\(h^{(d_C)}(\tuple v)=h^{(d_C)}(\tuple y)\), and template evaluation is
homomorphic.  Thus the same concrete parent context \(E\) accepts both parent
tuples, and Lemma~\ref{lem:shared-context} gives the conclusion.
\end{proof}

\begin{proposition}[Learner soundness]
\label{prop:learner-soundness}
Let \(L\) be \((f,h)\)-tuple-substitutable and let \(K\subseteq L\).  If
\([\tuple x]\Rightarrow^*\tuple u\) in \(\widehat G_0(K)\), then
\(\tuple u\equiv_L\tuple x\).  Consequently, \(L(\widehat G_0(K))\subseteq L\).
\end{proposition}

\begin{proof}
Induct on the derivation from \([\tuple x]\).  A terminal rule is immediate.  A
unit rule \([\tuple x]\to[\tuple y]\) was added only when \(h^{(d)}(\tuple x)=
h^{(d)}(\tuple y)\) and a concrete context \(E\) satisfies
\(E[\tuple x],E[\tuple y]\in K\subseteq L\).  Lemma~\ref{lem:shared-context}
gives \(\tuple x\equiv_L\tuple y\), and the induction hypothesis for
\([\tuple y]\) finishes the case.  For a binary rule
\([\tuple x]\to\rho([\tuple y],[\tuple z])\), the witness supplies a concrete
parent occurrence \((E,\tuple x)\) in \(K\), so \(E[\rho(\tuple y,\tuple z)]
=E[\tuple x]\in L\).  If the children derive \(\tuple u\) and \(\tuple v\), the
induction hypotheses give \(\tuple u\equiv_L\tuple y\) and
\(\tuple v\equiv_L\tuple z\).  Lemma~\ref{lem:witnessed-composition} yields
\(\rho(\tuple u,\tuple v)\equiv_L\tuple x\).

For the start symbol, any derivation begins with \(\widehat S\to[(w_0)]\) for
some \(w_0\in K\).  If it derives \(w\), then \((w)\equiv_L(w_0)\).  Since the
empty unary context \(\blank_1\) accepts \((w_0)\), it also accepts \((w)\), so
\(w\in L\).
\end{proof}

\subsection{Completeness and Exact Reconstruction}
\label{sec:completeness}

\begin{lemma}[Simulation of typed rules]
\label{lem:rule-simulation}
Assume \(\CS(\widetilde G_0)\subseteq K\subseteq L(G)\).  For each surviving
nonterminal \(X\), put \(\eta(X):=[\omega(X)]\).  Then every rule of
\(\widetilde G_0\) is simulated by \(\widehat G_0(K)\) under \(\eta\):
\begin{enumerate}[label=(\roman*),leftmargin=*]
\item If \(X\to(a)\), then \(\eta(X)\Rightarrow^*(a)\).
\item If \(X\to\rho(Y,Z)\), then
\(\eta(X)\Rightarrow^*\rho(\eta(Y),\eta(Z))\).
\item If \(\widetilde S\to X\), then \(\widehat S\to\eta(X)\) is a learner
start rule.
\end{enumerate}
\end{lemma}

\begin{proof}
For (i), the characteristic sample contains \(\chi(X)[\omega(X)]\) and
\(\chi(X)[(a)]\).  These two sample words use the same concrete context
\(\chi(X)\).  Also, by Proposition~\ref{prop:output-invariants},
\(h^{(1)}(\omega(X))=h(a)\).  Hence the learner has the unit rule
\([\omega(X)]\to[(a)]\), and it has the terminal rule \([(a)]\to(a)\).

For (ii), let
\[
  \tuple z_R:=\rho(\omega(Y),\omega(Z)).
\]
The characteristic sample contains both
\(\chi(X)[\omega(X)]\) and \(\chi(X)[\tuple z_R]\).  Hence
\(\omega(X)\) and \(\tuple z_R\) occur in the same concrete sentence context
\(\chi(X)\).  They have the same componentwise \(h\)-type because both belong
to the output-typed language \(L_X\).  Clause (iv) of
Definition~\ref{def:learner} therefore adds the unit rule
\[
  [\omega(X)]\to[\tuple z_R].
\]

The word \(\chi(X)[\tuple z_R]\) also supplies the parent occurrence used for
a binary witness.  By the definition of \(\tuple z_R\), every component of
\(\omega(Y)\) and \(\omega(Z)\) appears at the variable interval prescribed by
\(\rho\), with zero-length components represented by the fixed tie-order
convention.  The tuple values \(\omega(Y)\) and \(\omega(Z)\) are observed in
\(K\) because the characteristic sample separately contains
\(\chi(Y)[\omega(Y)]\) and \(\chi(Z)[\omega(Z)]\).  By
Remark~\ref{rem:binary-witness-locality}, these exposing occurrences need not
lie in the parent sample word.  Consequently the learner adds
\[
  [\tuple z_R]\to\rho([\omega(Y)],[\omega(Z)]).
\]
Combining the unit rule and the binary rule gives the required simulation.

For (iii), if \(\widetilde S\to X\) is a trimmed start rule, then by the
start-child convention \(\chi(X)=\blank_1\), so \(\omega(X)\in\CS(\widetilde
G_0)\subseteq K\).  Hence \(\widehat S\to[(\omega(X))]\) is a learner start
rule.
\end{proof}

\begin{proposition}[Completeness]
\label{prop:completeness}
If \(\CS(\widetilde G_0)\subseteq K\subseteq L(G)\), then
\(L(G)\subseteq L(\widehat G_0(K))\).
\end{proposition}

\begin{proof}
We prove by induction on derivation height that, for every surviving typed
nonterminal \(X\), every tuple derived from \(X\) in \(\widetilde G_0\) is
derivable from \(\eta(X)=[\omega(X)]\) in the learner grammar.  For a terminal
rule this is Lemma~\ref{lem:rule-simulation}(i).  In the binary case
\(X\to\rho(Y,Z)\), apply the induction hypotheses to the two child derivations
and prepend the finite simulation segment supplied by
Lemma~\ref{lem:rule-simulation}(ii).  This derives the same parent tuple from
\(\eta(X)\).

By Proposition~\ref{prop:output-invariants}, every successful \(G\)-derivation
lifts to a successful derivation in \(\widetilde G_0\).  Its typed start rule is
simulated by Lemma~\ref{lem:rule-simulation}(iii), and the induction just proved
simulates the remaining tree.  Hence every word of \(L(G)\) is generated by
\(\widehat G_0(K)\).
\end{proof}

\begin{theorem}[Exact reconstruction]
\label{thm:exact-reconstruction}
Fix \(f\) and an explicit finite homomorphism \(h\).  Let \(G\) be a reduced
working binary linear nondeleting MCFG of fan-out at most \(f\), and let
\(L=L(G)\) be \((f,h)\)-tuple-substitutable.  If \(\widetilde G_0\) is the
trimmed output-type refinement of \(G\), then every finite sample \(K\) with
\(\CS(\widetilde G_0)\subseteq K\subseteq L\) satisfies
\(L(\widehat G_0(K))=L\).
\end{theorem}

\begin{proof}
Soundness is Proposition~\ref{prop:learner-soundness}.  Completeness is
Proposition~\ref{prop:completeness}.
\end{proof}

\begin{remark}[Language-level exactness]
\label{rem:language-level-exactness}
The reconstruction is exact at the language level.  The learner need not
recover the witnessing nonterminal names, its derivation trees, or an
isomorphic copy of the target presentation.
\end{remark}

\begin{corollary}[Identification in the limit]
\label{cor:gold}
For fixed \(f\) and fixed explicit \(h\), the class \(\Cmcf{f}{h}\) is
identifiable in the limit from positive data by the canonical learner
\(K\mapsto\widehat G_0(K)\).  Recall that \(\Cmcf{f}{h}\) is defined through
the existence of a reduced working binary linear nondeleting MCFG presentation,
not through arbitrary MCFG presentations.
\end{corollary}

\begin{proof}
For each target \(L\in\Cmcf{f}{h}\), choose a witnessing grammar \(G\) and its
finite characteristic sample \(\CS(\widetilde G_0)\).  Apply
Lemma~\ref{lem:sufficient-to-gold} to Theorem~\ref{thm:exact-reconstruction}.
\end{proof}

\begin{remark}[Possible extension to bounded rule rank]
\label{rem:bounded-rule-rank}
The binary-rule convention is used throughout the formal development.  The
same concrete-witness strategy appears to extend to every fixed rule-rank
bound \(r\): one would segment a parent tuple into the components of at most
\(r\) observed child tuples and replace the children one at a time in the
soundness argument.  A complete formulation would require corresponding
\(r\)-ary witness definitions, enumeration bounds, and simulation lemmas.  We
leave that extension to future work and make no rank-\(r\) theorem here.
\end{remark}



\section{Hypothesis Construction and Exposure Size}
\label{sec:poly-time}

\begin{theorem}[Fixed-parameter polynomial construction]
\label{thm:poly}
For fixed \(f\) and fixed explicit \(h\), the grammar \(\widehat G_0(K)\) is
constructible in time \(\|K\|_+^{O(f)}\), with constants depending on the
explicit monoid.  In particular, for each fixed \(f\) and \(h\), this is a
polynomial in \(\|K\|_+\), including output size; the degree of that polynomial
grows with \(f\), so the bound is fixed-parameter rather than uniform in
\(f\).
\end{theorem}

\begin{proof}
Let \(n=\|K\|_+\).  For fixed \(d\le f\), an arity-\(d\) tuple occurrence in a
word of length at most \(n\) is determined by \(O(d)\) cuts or interval endpoints,
together with an ordering datum depending only on \(d\).  Hence there are
\(n^{O(d)}\) arity-\(d\) occurrences, and summing over \(d\le f\) and over the
words in \(K\) gives at most \(n^{O(f)}\) observed tuples.

Terminal and start rules are immediate.  Unit rules are found by comparing
pairs of observed tuples of the same arity, checking equality of componentwise
\(h\)-types, and enumerating observed concrete occurrences to decide whether a
shared concrete context fills both tuples to sample words.  All of these
operations are bounded by \(n^{O(f)}\) because \(M\) and \(f\) are fixed.

For binary witness rules, enumerate a parent occurrence, two observed child
tuple values of arities at most \(f\), and placements of their components in a
segmentation of the parent components into terminal gaps and at most \(2f\)
variable intervals.  Since \(f\) is fixed, the
number of endpoint choices, cut choices, and tie-order choices is
\(n^{O(f)}\).  Each candidate segmentation is checked by direct string
comparison and then output as one rule.  Hence the number of generated rules
and the time needed to generate them are \(n^{O(f)}\), including output size.

Unit rules are part of the learner grammar.  If a unit-free working grammar is
required, compute the unit closure over the polynomial-size nonterminal set and
replace unit paths by the corresponding terminal and binary rules.  This causes
only polynomial blow-up for fixed \(f\) and \(h\).
\end{proof}

\subsection{Exposure size and the polynomial-data boundary}
\label{subsec:exposure-boundary}

\begin{definition}[Exposure size]
\label{def:exposure-size}
Let \(G\) be a witnessing target presentation and let \(\widetilde G_0\) be its
trimmed output-type refinement.  Define \(B_{\mathrm{exp}}(G,h)\) as the
maximum positive size of the characteristic-sample words
\[
  \chi(X)[\omega(X)],\qquad
  \chi(X)[(a)],\qquad
  \chi(X)[\rho(\omega(Y),\omega(Z))]
\]
used for surviving typed nonterminals, terminal rules, and binary rules of
\(\widetilde G_0\).
\end{definition}

\begin{proposition}[Characteristic-sample size under an exposure bound]
\label{prop:cs-exposure-bound}
Let \(N_{\mathrm{NT}}\) be the number of surviving typed nonterminals of
\(\widetilde G_0\), and let \(N_{\mathrm{rule}}\) be the number of surviving
non-start typed rules.  Then
\[
  |\CS(\widetilde G_0)|
  \le N_{\mathrm{NT}}+N_{\mathrm{rule}}
\]
and
\[
  \|\CS(\widetilde G_0)\|_+
  \le
  (N_{\mathrm{NT}}+N_{\mathrm{rule}})
  B_{\mathrm{exp}}(G,h).
\]
Consequently, every subclass for which \(N_{\mathrm{NT}}\),
\(N_{\mathrm{rule}}\), and \(B_{\mathrm{exp}}(G,h)\) are polynomially bounded
in the chosen witnessing-presentation size has polynomial-size characteristic
samples.
\end{proposition}

\begin{proof}
Definition~\ref{def:cs} contributes at most one anchor word for every surviving
typed nonterminal and at most one filled word for every surviving non-start
typed rule; duplicate words can only reduce the cardinality.  Every contributed
word has positive size at most \(B_{\mathrm{exp}}(G,h)\), which gives the two
bounds.
\end{proof}

The full output-type refinement has at most
\[
  \sum_{A\in V\setminus\{S\}} |M|^{\mu(A)}
\]
typed nonterminals, and a binary rule \(A\to\rho(B,C)\) gives at most
\(|M|^{\mu(B)+\mu(C)}\) typed rules.  Thus, for fixed \(f\) and fixed \(h\),
the combinatorial size of the refinement is bounded by a constant multiple of
the witnessing-grammar size.  The uncontrolled quantity in the full class is
the length of the exposing derivations and hence the exposure bound.

This proposition does not establish a polynomial exposure bound for the whole
class \(\Cmcf{f}{h}\).  The unconditional result of this paper is finite exact
reconstruction together with polynomial-time construction of the hypothesis
from a given finite sample.  A polynomial time-and-data theorem is obtained
only for subclasses with an additional polynomial exposure bound.


\section{The Fixed Observation Parameter}
\label{sec:fixed-observation}

The fixed morphism \(h\) is part of the learning problem in
Corollary~\ref{cor:gold}; the learner is not asked to discover a suitable
finite congruence from positive data.  This section shows that the assumption is
not merely a presentation convenience.  A fixed monoid-size bound can be
compiled into one product observation, whereas the unbounded union over all
finite observations is not identifiable already in fan-out one.

\begin{definition}[No-advice union class]
\label{def:no-advice-union}
For a fixed alphabet and a fixed fan-out bound \(f\), define
\[
  \Cmcf{f}{\bullet}:=\bigcup_h \Cmcf{f}{h},
\]
where the union ranges over all explicit finite monoid homomorphisms
\(h\colon\Sigma^*\to M\).  A learner for this class receives positive data only;
it is not given which morphism witnesses membership of the target.
\end{definition}

\begin{lemma}[Ascending-chain obstruction]
\label{lem:ascending-chain-obstruction}
Let a language class \(\mathcal C\) contain languages
\[
  L_1\subsetneq L_2\subsetneq\cdots
\]
and their union \(L_\infty=\bigcup_{k\ge1}L_k\).  Then \(\mathcal C\) is not
identifiable in the limit from positive data.
\end{lemma}

\begin{proof}
Assume that a learner \(\mathcal A\) identifies every language in
\(\mathcal C\).  We construct a text for \(L_\infty\) on which
\(\mathcal A\) outputs infinitely many different hypothesis languages.
Fix an enumeration \(v_1,v_2,\ldots\) of \(L_\infty\).  Suppose a finite prefix
\(\sigma_s\) has been constructed and its content is contained in some
\(L_{k_s}\).  Continue \(\sigma_s\) with elements of a text for \(L_{k_s}\).
The resulting infinite continuation is a text for \(L_{k_s}\), so at some
finite extension the learner must output a hypothesis whose language is
\(L_{k_s}\).  Then append an element of
\(L_{k_s+1}\setminus L_{k_s}\), and append \(v_s\) if it has not yet appeared.
The new finite content is contained in some later \(L_{k_{s+1}}\), and the
construction can be repeated.

Every \(v_s\) eventually appears, so the resulting sequence is a text for
\(L_\infty\).  At the selected stages, however, the learner outputs the
strictly increasing approximant languages \(L_{k_s}\), and therefore cannot
converge semantically on this text.  This contradicts identification of
\(L_\infty\).
\end{proof}

Let
\[
  L_*:=\{a^n b^n\mid n\ge1\},
  \qquad
  L_k:=\{a^n b^n\mid 1\le n\le k\}.
\]
All these languages have fan-out-one working binary linear nondeleting MCFG
presentations.  We first put the limit language into one fixed observation
fiber, for every arity bound.

\begin{lemma}[The limit language is substitutable for a coarse envelope morphism]
\label{lem:limit-ab-substitutable}
Let \(h_*\) be the transition morphism of the standard DFA for \(a^*b^*\).  Then
\(L_*\) is \((f,h_*)\)-tuple-substitutable for every \(f\ge1\).
\end{lemma}

\begin{proof}
We prove the condition for every arity \(d\).  Suppose tuples
\(\tuple x,\tuple y\in(\Sigma^*)^d\) have the same componentwise \(h_*\)-type
and share an accepting sentence context \(E\).  Thus
\(E[\tuple x]=a^n b^n\) and \(E[\tuple y]=a^m b^m\) for some \(m,n\ge1\).
Let \(A_x,B_x\) be the total numbers of \(a\)- and \(b\)-letters occurring in
the components of \(\tuple x\), and define \(A_y,B_y\) similarly.  The outside
terminal contribution of \(E\) is the same in the two fillings, so comparison
of the two accepted words gives
\[
  A_y-A_x=m-n=B_y-B_x .
\]

Now let \(F\) be any accepting sentence context for \(\tuple x\), say
\(F[\tuple x]=a^N b^N\).  Replacing the components one at a time by components
with the same \(h_*\)-value preserves the transition morphism of the whole
filled sentence.  Hence \(F[\tuple y]\in a^*b^*\).  Its total number of
\(a\)-letters is \(N+(A_y-A_x)\), and its total number of \(b\)-letters is
\(N+(B_y-B_x)\), which are equal by the displayed identity.  This common value
cannot be zero: if \(F[\tuple y]=\eps\), then all terminal parts of \(F\) and
all components of \(\tuple y\) are empty; since no nonempty word acts as the
identity transformation of the standard \(a^*b^*\) zone automaton, equality of
componentwise \(h_*\)-types forces all components of \(\tuple x\) to be empty,
contrary to \(F[\tuple x]\in L_*\).  Therefore \(F[\tuple y]\in L_*\).  The
reverse inclusion is symmetric, so the tuple distributions are equal.
\end{proof}

\begin{lemma}[Finite approximants can be made substitutable]
\label{lem:finite-approximants-substitutable}
For every \(k\ge1\) and every \(f\ge1\), there exists an explicit finite monoid
morphism \(h_k\) such that \(L_k\) is \((f,h_k)\)-tuple-substitutable.  Moreover,
if a morphism \(h\) makes \(L_k\) \((1,h)\)-tuple-substitutable, then the values
\(h(a^i b^i)\), \(1\le i\le k\), are pairwise distinct.  In particular
\(|\operatorname{im}(h)|\ge k\).
\end{lemma}

\begin{proof}
For existence, take \(h_k\) to be the syntactic morphism of the regular language
\(L_k\).  If \(h_k(x_i)=h_k(y_i)\) for each component of two arity-\(d\) tuples
\(\tuple x,\tuple y\), then replacing the components one at a time preserves
membership in \(L_k\) in every surrounding sentence context, because each
replacement is by syntactically equivalent strings with respect to \(L_k\).
Hence \(\D_{L_k}^{(d)}(\tuple x)=\D_{L_k}^{(d)}(\tuple y)\) for every \(d\), and
the \((f,h_k)\)-substitutability implication is immediate.

For the lower bound, fix \(1\le i<j\le k\).  The words
\(x_i=a^i b^i\) and \(x_j=a^j b^j\) share the accepting context
\((\eps,\eps)\), since both belong to \(L_k\).  Their distributions differ:
\((a^{k-i},b^{k-i})\) accepts \(x_i\), producing \(a^k b^k\in L_k\), but it
sends \(x_j\) to \(a^{k-i+j}b^{k-i+j}\notin L_k\).  Therefore, if
\(h(x_i)=h(x_j)\), the \((1,h)\)-substitutability implication would fail.  All
\(h(a^i b^i)\), \(1\le i\le k\), must be distinct.
\end{proof}

\begin{proposition}[The separating chain cannot lie in one fixed fiber]
\label{prop:separating-chain-not-one-fiber}
For any fixed finite morphism \(h\), only finitely many of the languages
\(L_k\) can be \((1,h)\)-tuple-substitutable.  More precisely, if
\(L_k\) is \((1,h)\)-tuple-substitutable, then \(k\le |\operatorname{im}(h)|\).
\end{proposition}

\begin{proof}
This is the lower bound in Lemma~\ref{lem:finite-approximants-substitutable}.
For fixed \(h\), the image \(\operatorname{im}(h)\) is finite, so only those
\(k\le |\operatorname{im}(h)|\) can occur.  Thus the superfinite chain used
below necessarily moves through increasingly refined observation fibers; it is
not a contradiction to the fixed-\(h\) learnability theorem.
\end{proof}

\begin{theorem}[No-advice non-identifiability]
\label{thm:no-advice-nonidentifiability}
For every fixed \(f\ge1\), the no-advice class \(\Cmcf{f}{\bullet}\) is not
identifiable in the limit from positive data.
\end{theorem}

\begin{proof}
By Lemma~\ref{lem:limit-ab-substitutable},
\(L_*\in\Cmcf{f}{h_*}\), and by
Lemma~\ref{lem:finite-approximants-substitutable}, every
\(L_k\in\Cmcf{f}{h_k}\) for some explicit finite morphism \(h_k\).  Thus the
no-advice union contains \(L_1\subsetneq L_2\subsetneq\cdots\) and their union
\(L_*\).  Lemma~\ref{lem:ascending-chain-obstruction} applies.
\end{proof}

The theorem should be read together with Corollary~\ref{cor:gold}.  A
fixed finite observation morphism selects one learnable fiber.  If the fiber is
not supplied and the target may use any finite morphism, the union class
contains a superfinite chain and positive data alone cannot determine the
correct finite observation interface.

\begin{definition}[Bounded-size observation union]
\label{def:bounded-observation-union}
Fix the finite alphabet \(\Sigma\), a fan-out bound \(f\ge1\), and \(k\ge1\).
Define
\[
  \Cmcf{f}{\le k}:=\bigcup_h \Cmcf{f}{h},
\]
where the union ranges over all explicit finite monoid homomorphisms
\(h:\Sigma^*\to M\) with \(|M|\le k\).
\end{definition}

\begin{lemma}[Isomorphism invariance of the observation]
\label{lem:iso-invariance}
If \(\varphi:M\to M'\) is an injective monoid homomorphism, then a language is
\((f,h)\)-tuple-substitutable iff it is \((f,\varphi\circ h)\)-tuple-substitutable.
Consequently \(\Cmcf{f}{h}=\Cmcf{f}{\varphi\circ h}\).
\end{lemma}

\begin{proof}
Definition~\ref{def:tuple-substitutable} uses \(h\) only through equalities of
componentwise values.  Since \(\varphi\) is injective,
\(h(u)=h(v)\) iff \(\varphi(h(u))=\varphi(h(v))\).  The grammar-theoretic part of
membership is unchanged.
\end{proof}

\begin{theorem}[Bounded observation size restores identifiability]
\label{thm:bounded-union-identifiable}
For every fixed finite alphabet \(\Sigma\), fan-out bound \(f\ge1\), and
\(k\ge1\), there is an explicit finite monoid homomorphism
\(H=H_{\Sigma,k}:\Sigma^*\to M_H\) such that
\[
  \Cmcf{f}{\le k}\subseteq\Cmcf{f}{H}.
\]
Consequently the canonical learner with parameters \((f,H)\) identifies
\(\Cmcf{f}{\le k}\) in the limit from positive data.
\end{theorem}

\begin{proof}
Up to isomorphism there are only finitely many monoids of cardinality at most
\(k\); fix explicit representatives \(M_1,\ldots,M_t\).  Since \(\Sigma^*\) is
free, a homomorphism \(\Sigma^*\to M_j\) is uniquely determined by an arbitrary
map \(\Sigma\to M_j\), so there are at most \(t k^{|\Sigma|}\) homomorphisms into
the representatives.  Enumerate them as
\(h_i:\Sigma^*\to M_{j_i}\), \(1\le i\le r\), and define
\[
  M_H:=M_{j_1}\times\cdots\times M_{j_r},
  \qquad
  H:=(h_1,\ldots,h_r):\Sigma^*\to M_H,
\]
with componentwise multiplication.  The multiplication table of \(M_H\) and the
letter values \(H(a)\), \(a\in\Sigma\), are computable from the factor tables, so
\(H\) is explicit.

Let \(h:\Sigma^*\to M\) be any explicit homomorphism with \(|M|\le k\).  Choose
an isomorphism \(\varphi:M\to M_j\) onto a representative.  Then
\(\varphi\circ h\) appears in the enumeration, say \(\varphi\circ h=h_i\).  The
homomorphism
\[
  \pi:=\varphi^{-1}\circ\mathrm{pr}_i:M_H\to M
\]
satisfies \(\pi\circ H=h\).  Hence \(h\preceq H\), and
Proposition~\ref{prop:h-refinement-monotonicity} gives
\(\Cmcf{f}{h}\subseteq\Cmcf{f}{H}\).  Taking the union over all \(|M|\le k\)
yields \(\Cmcf{f}{\le k}\subseteq\Cmcf{f}{H}\).  By
Corollary~\ref{cor:gold}, the canonical learner for the fixed morphism \(H\)
identifies \(\Cmcf{f}{H}\).  Identification in the limit is a universal
statement over targets and texts, so the same learner identifies the subclass
\(\Cmcf{f}{\le k}\).
\end{proof}

\begin{remark}[Existence rather than practical size]
\label{rem:universal-product-size}
The universal product morphism \(H_{\Sigma,k}\) may be extremely large.  The
theorem is an identifiability result for each fixed \((\Sigma,f,k)\), not a
claim that this product is a practical representation or that its size is
polynomial in \(k\).  Its role is to show that the failure of the no-advice
union begins only when the observation size is unbounded.
\end{remark}

\begin{remark}[Bounded versus unbounded observation advice]
\label{rem:bounded-vs-unbounded-union}
Theorem~\ref{thm:bounded-union-identifiable} does not contradict
Theorem~\ref{thm:no-advice-nonidentifiability}.  It locates the failure at the
unboundedness of the observation monoid.  By
Lemma~\ref{lem:finite-approximants-substitutable}, the finite approximant
\(L_k=\{a^n b^n\mid 1\le n\le k\}\) is \((1,h)\)-tuple-substitutable only if
\(|\operatorname{im}(h)|\ge k\).  Thus the superfinite chain used in the
no-advice theorem cannot lie in any fixed bounded-size slice.  Since
\(\Cmcf{f}{\bullet}=\bigcup_{k\ge1}\Cmcf{f}{\le k}\), positive-data
identifiability holds on every bounded slice but is not preserved under the
increasing union.  This is an identifiability statement, not an efficiency
claim uniform in \(k\) or \(\Sigma\): for fixed \((\Sigma,k)\), the product
morphism \(H_{\Sigma,k}\) is fixed, but its size may be enormous.
\end{remark}

\subsection{Infinite member kernels}
\label{subsec:member-kernel-exclusion}

For a language \(L\subseteq\Sigma^*\), define its arity-one member kernel by
\[
  \mathrm{MK}_1(L):=\{\D_L^{(1)}(w)\mid w\in L\}.
\]
This is not a finite-kernel condition for all residual distributions; it only
records distributions of strings that are themselves members of \(L\).

\begin{theorem}[Exclusion by infinite member kernel]
\label{thm:member-kernel-exclusion}
Let \(L\subseteq\Sigma^*\).  If \(\mathrm{MK}_1(L)\) is infinite, then
\[
  L\notin \bigcup_h \Cmcf{f}{h}
\]
for every \(f\ge1\), where the union ranges over all explicit finite monoid
homomorphisms \(h:\Sigma^*\to M\).
\end{theorem}

\begin{proof}
Suppose, toward contradiction, that \(L\in\Cmcf{f}{h}\) for some explicit finite
monoid morphism \(h:\Sigma^*\to M\).  Let \(w,w'\in L\) satisfy \(h(w)=h(w')\).
The identity sentence context \(\blank_1\) is accepting for both \(w\) and
\(w'\).  Since \(f\ge1\) and \(L\) is \((f,h)\)-tuple-substitutable, the
arity-one implication gives
\[
  \D_L^{(1)}(w)=\D_L^{(1)}(w').
\]
Thus the distribution \(\D_L^{(1)}(w)\), for member strings \(w\in L\), is
determined by the finite value \(h(w)\in M\).  Hence at most \(|M|\) different
member distributions can occur, contradicting the infinitude of
\(\mathrm{MK}_1(L)\).
\end{proof}

\begin{corollary}[The slope union is outside every finite observation]
\label{cor:slope-union-outside-every-finite-observation}
For every \(f\ge1\),
\[
  L_{\mathrm{slope}}
  =
  \{a^n b^n\mid n\ge1\}
  \cup
  \{a^n b^{2n}\mid n\ge1\}
\]
satisfies
\[
  L_{\mathrm{slope}}\notin \bigcup_h \Cmcf{f}{h}.
\]
\end{corollary}

The slope union is context-free (it is the union of two context-free
languages), so this obstruction is not caused by high fan-out or by copying
power.

\begin{proof}
By Theorem~\ref{thm:member-kernel-exclusion}, it suffices to show that the
arity-one member kernel is infinite.  For \(n\ge1\), put
\[
  w_n=a^n b^n\in L_{\mathrm{slope}},
  \qquad
  C_n=\blank_1 b^n .
\]
Then \(C_n[w_n]=a^n b^{2n}\in L_{\mathrm{slope}}\).  If \(m\ne n\), then
\(C_n[w_m]=a^m b^{m+n}\).  This word is not in \(\{a^r b^r\mid r\ge1\}\), since
that would require \(m+n=m\).  It is not in \(\{a^r b^{2r}\mid r\ge1\}\) either,
since that would require \(m+n=2m\), hence \(n=m\), contrary to assumption.
Therefore
\[
  C_n\in\D_{L_{\mathrm{slope}}}^{(1)}(w_n)
  \quad\text{but}\quad
  C_n\notin\D_{L_{\mathrm{slope}}}^{(1)}(w_m)
  \qquad(m\ne n).
\]
The member distributions are pairwise distinct, so
Theorem~\ref{thm:member-kernel-exclusion} applies.
\end{proof}

\begin{proposition}[The copy language is outside every finite observation]
\label{prop:copy-language-outside-every-finite-observation}
Let
\[
  L_{\mathrm{copy}}=
  \{ww\mid w\in\{a,b\}^+\}.
\]
Then
\[
  L_{\mathrm{copy}}\notin\bigcup_h \Cmcf{f}{h}
\]
for every \(f\ge1\).  This holds even though \(L_{\mathrm{copy}}\) is a standard
fan-out-two linear multiple context-free language.
\end{proposition}

\begin{proof}
The fan-out-two upper bound is the usual MCFG copy construction; for example,
with rank-one linear rules one may use
\(A\to(a x^1,a x^2)(A)\), \(A\to(b x^1,b x^2)(A)\),
\(A\to(a,a)\), and \(A\to(b,b)\), followed by a top concatenation rule.  This
upper bound is only background here; the exclusion below is independent of any
presentation.

We prove that the arity-one member kernel is infinite.  For \(n\ge1\), put
\[
  x_n=a^n b,
  \qquad
  z_n=x_nx_n\in L_{\mathrm{copy}},
  \qquad
  C_n=\blank_1 z_n .
\]
Then \(C_n[z_n]=z_nz_n\in L_{\mathrm{copy}}\).  We claim that if \(m\ne n\), then
\[
  C_n[z_m]=z_mz_n=x_mx_mx_nx_n
\]
is not a square.  Write
\[
  x_mx_mx_nx_n=a^m b\,a^m b\,a^n b\,a^n b,
\]
and let \(H=m+n+2\), half of its total length.  In any square, letters at
positions separated by \(H\) agree.  The four occurrences of \(b\) are at
positions
\[
  m+1,
  \quad 2m+2,
  \quad 2m+n+3,
  \quad 2m+2n+4 .
\]
If \(m<n\), the second position \(2m+2\) lies in the first half, while
\[
  (2m+2)+H=3m+n+4
\]
lies strictly between the third and fourth occurrences of \(b\), where the
letter is \(a\).  This contradicts the square condition.  If \(m>n\), then the
second position lies in the second half, and shifting it back by \(H\) gives
\[
  (2m+2)-H=m-n,
\]
which lies before the first occurrence of \(b\), again at an \(a\).  This also
contradicts the square condition.

Thus
\[
  C_n\in\D_{L_{\mathrm{copy}}}^{(1)}(z_n)
  \quad\text{but}\quad
  C_n\notin\D_{L_{\mathrm{copy}}}^{(1)}(z_m)
  \qquad(m\ne n).
\]
The member distributions \(\D_{L_{\mathrm{copy}}}^{(1)}(z_n)\) are pairwise
distinct.  Theorem~\ref{thm:member-kernel-exclusion} excludes
\(L_{\mathrm{copy}}\) from every finite-observation class.
\end{proof}

\begin{remark}
The slope-union example shows that the obstruction is not merely a poor choice
of regular envelope: no finite observation morphism can make that union safe.
The copy-language example shows a different boundary: even a standard fan-out-two
MCFL can fail every finite-observation promise.
\end{remark}



\section{Comparison with Distributional Learning}
\label{sec:comparison}

The fixed-observation condition is best viewed as a typed version of
distributional substitutability.  It weakens the untyped trigger by requiring a
componentwise finite-monoid type match before shared contexts force equality of
distributions.

\paragraph{Three differences from multidimensional substitutability.}
The present condition is close in spirit to Yoshinaka's multidimensional
substitutability, but it is not intended as a direct replacement of that
condition.  There are three relevant differences.

First, the substitution trigger is typed.  In the untyped setting, a shared
accepting multidimensional context already forces equality of distributions.
Here a shared accepting named sentence context forces equality only after the
componentwise finite observation \(h^{(d)}\) agrees.  Thus the finite morphism
acts as a regular or finite-state guard on distributional substitution.

Second, the contexts used here are sentence-level named-hole contexts rather
than grammar-internal multidimensional contexts.  They record how a tuple is
exposed in complete strings, with named components allowed to appear in any
sentence order.  This is deliberately weaker than storing a full derivational
interface: the learner observes concrete occurrences in positive strings and
uses the characteristic sample to supply the missing exposing contexts.

Third, the result is advice-relative.  The finite morphism \(h\) is fixed before
learning starts and is not inferred from the text.  Consequently, the theorem
identifies a fixed observation fiber \(\Cmcf{f}{h}\), not the full class of
MCFLs and not the union over all possible finite observations.

\begin{center}
\begin{tabular}{|p{0.28\textwidth}|p{0.30\textwidth}|p{0.30\textwidth}|}
\hline
 & Multidimensional substitutability & Fixed-observation tuple substitutability \\
\hline
Trigger & Shared multidimensional context & Shared named sentence context plus equal \(h\)-type \\
\hline
Type information & None, or implicit in the context system & Explicit finite monoid observation fixed in advance \\
\hline
Learning target & A distributionally restricted MCFL class & A fixed advice-relative fiber \(\Cmcf{f}{h}\) \\
\hline
Main limitation & Context condition controls substitution directly & The observation \(h\) is not learned from positive data \\
\hline
\end{tabular}
\end{center}

The examples in Section~\ref{sec:running-examples} illustrate why the finite
guard is useful.  In block-synchronization languages, a shared sentence context
alone is too coarse: tuples may occur in the same outside position while
crossing different regular zones.  The transition morphism of the regular
envelope prevents such unsafe substitutions, while the equality constraints are
still learned distributionally from positive data.

Let \(\Yclass{f}^{\mathrm{nam}}\) denote the class of fan-out-at-most-\(f\) MCFLs
satisfying the untyped named-sentence substitutability condition: for every
\(1\le d\le f\), whenever two arity-\(d\) tuples share an accepting named
sentence context, their accepting named sentence-context distributions are equal.

\begin{proposition}[Untyped named substitutability implies every fixed observation]
\label{prop:untyped-named-implies-fixed-h}
For every explicit finite morphism \(h\) and every \(f\ge1\),
\[
  \Yclass{f}^{\mathrm{nam}}\subseteq \Cmcf{f}{h}.
\]
\end{proposition}

\begin{proof}
The fixed-\(h\) implication has the same conclusion as the untyped named
condition and only adds the extra antecedent
\(h^{(d)}(\tuple x)=h^{(d)}(\tuple y)\).  Thus every untyped named-substitutable
language is automatically \((f,h)\)-tuple-substitutable.  The same fan-out
presentation witnesses the grammar-theoretic part of membership.
\end{proof}

The inclusion can be strict.  Let \(L_3=\{a^n b^n c^n\mid n\ge1\}\), and let
\(h_R\) be the transition morphism of the standard DFA for \(R=a^*b^*c^*\).
The language \(L_3\) has a fan-out-two working presentation by
Proposition~\ref{prop:l3-running-example}.  It is also \((f,h_R)\)-tuple-
substitutable for every \(f\ge2\).  Indeed, a shared accepting context fixes a
common shift in the three letter counts, while equality of \(h_R\)-types keeps
every later filling inside the block zone \(a^*b^*c^*\); hence every later
accepting context for one tuple is accepting for the other, and conversely.

On the other hand, \(L_3\) is not untyped named-substitutable.  The arity-two
tuples
\[
  \tuple x=(a,c),
  \qquad
  \tuple y=(a^2b,c^2)
\]
share the accepting context \(E=\blank_1 b\blank_2\), since
\(E[\tuple x]=abc\) and \(E[\tuple y]=a^2b^2c^2\).  But for
\(F=\blank_1 abb\blank_2 c\), one has
\(F[\tuple x]=a^2b^2c^2\in L_3\), whereas
\(F[\tuple y]=aababbccc\notin L_3\).  Thus the typed fixed-observation condition
strictly extends the corresponding untyped named condition.

The comparison with Yoshinaka's ordered multidimensional contexts is not a
literal inclusion statement, because ordered holes and named sentence holes test
different context geometries.  The present formalism deliberately uses named
holes, since MCFG components may be exposed in a surrounding sentence in a
permuted order.  The fixed morphism then supplies a finite type discipline on
those exposed components rather than an untyped global equality of all shared
context distributions.



\section{Conclusion}
\label{sec:conclusion}

For every fixed fan-out bound \(f\) and fixed explicit finite monoid observation
\(h\), the class \(\Cmcf{f}{h}\) is identifiable in the limit from positive
data by a canonical tuple-enumerating learner.  The construction stores
componentwise output types and observed tuple values; component placement is
recovered from concrete witnessed segmentations, and substitution is guarded
by an actually shared sentence context.  The reconstruction is exact at the
language level, while the observation morphism and the semantic
substitutability promise remain external assumptions.

The observation parameter has a sharp qualitative boundary.  A fixed fiber is
learnable; the union of all fibers with a fixed monoid-size bound is still
learnable through one finite product morphism; but the unbounded union over all
finite observations is not identifiable from positive data.  Independently,
the infinite member-kernel criterion excludes languages whose member
context-distributions cannot be compressed into finitely many observation
values.

The resulting region is nontrivial but proper.  It contains standard
three-block, parallel, and cross-serial synchronization examples, while the
context-free slope union and the fan-out-two copy language lie outside every
finite-observation fiber.  Fixed finite-monoid advice therefore isolates a
genuine learnable region inside the multiple context-free languages rather
than a disguised presentation of the full class.


\appendix
\section{Semantic Non-Effectivity of the Promise}
\label{app:semantic-promise}

The identification theorem is conditional on the semantic
\((f,h)\)-tuple-substitutability promise.  This appendix records why the promise
is not treated as an effectively checkable syntactic side condition.

\begin{theorem}[Undecidability for arbitrary fixed finite observations]
\label{thm:promise-undecidable}
Let \(\Sigma\) be a fixed finite alphabet with \(|\Sigma|\ge2\), and let
\(h:\Sigma^*\to M\) be any fixed explicit homomorphism into a finite monoid.
Then it is undecidable, given a context-free grammar over \(\Sigma\), whether
its language is \((1,h)\)-tuple-substitutable.  The same holds for
\((f,h)\)-tuple-substitutability for every fixed \(f\ge1\).  Moreover, the set
of positive instances is co-r.e.-complete (equivalently,
\(\Pi^0_1\)-complete under standard effective encodings).
\end{theorem}

\begin{proof}
Universality for context-free grammars is \(\Pi^0_1\)-complete under
standard effective encodings \cite{DomaratzkiOkhotinShallit2007}.  We first
record that the same classification holds over the fixed alphabet \(\Sigma\).
Given a CFG \(G\) over an arbitrary finite alphabet \(\Gamma\), choose distinct
letters \(a,b\in\Sigma\) and an injective fixed-length coding
\(c:\Gamma\to\{a,b\}^m\), extended homomorphically to \(\Gamma^*\).  The
language \(C:=c(\Gamma)^*\) is regular, and one can effectively construct a CFG
over \(\Sigma\) for
\[
  U_G:=(\Sigma^*\setminus C)\cup c(L(G)).
\]
Since \(c\) is injective,
\[
  U_G=\Sigma^*
  \quad\Longleftrightarrow\quad
  L(G)=\Gamma^*.
\]
Thus universality is \(\Pi^0_1\)-hard over every fixed alphabet of cardinality
at least two.  It belongs to \(\Pi^0_1\), since nonuniversality is recursively
enumerable by enumerating words and deciding CFG membership.  Hence fixed-
alphabet universality is \(\Pi^0_1\)-complete.  The nonunary restriction is
essential: Hopcroft's characterization of equality with a fixed regular
language implies that CFG universality is decidable over a unary alphabet and
undecidable over every fixed nonunary alphabet \cite{Hopcroft1969}.

We now reduce fixed-alphabet CFG universality to the semantic promise.  Let
\(G_0\) be a CFG over \(\Sigma\).  Since \(M\) is finite, choose an integer
\(N\) with \(|\Sigma|^N>|M|\).  Among the words of
length \(N\) there are distinct words \(X,Y\in\Sigma^N\) with
\(h(X)=h(Y)\).  The equal length and distinctness imply
\(X\Sigma^*\cap Y\Sigma^*=\emptyset\).

Fix a nonempty word \(\beta\in\Sigma^+\).  Given a grammar \(G_0\) over
\(\Sigma\), define
\[
  L'(G_0):=(\Sigma^*\setminus \beta\Sigma^*)\cup \beta L(G_0).
\]
Then \(L'(G_0)=\Sigma^*\) if \(L(G_0)=\Sigma^*\), while if
\(z\notin L(G_0)\) then \(\beta z\notin L'(G_0)\).  Also
\(\eps\in L'(G_0)\).

Let \(H_0=h(\Sigma^*)\) be the image submonoid and set
\[
  S:=h(Y)H_0,
  \qquad K_{\mathrm{pos}}:=h^{-1}(S).
\]
Finally put
\[
  K(G_0):=\bigl(K_{\mathrm{pos}}\setminus Y\Sigma^*\bigr)\cup YL'(G_0).
\]
The language \(K(G_0)\) is effectively context-free: \(K_{\mathrm{pos}}\) and
\(Y\Sigma^*\) are regular, and \(YL'(G_0)\) is context-free.

If \(L(G_0)=\Sigma^*\), then \(L'(G_0)=\Sigma^*\).  Since
\(Y\Sigma^*\subseteq K_{\mathrm{pos}}\), we obtain
\[
  K(G_0)=K_{\mathrm{pos}}=h^{-1}(S).
\]
Thus membership in \(K(G_0)\) is determined by the \(h\)-value of the whole
word.  If two tuples have the same componentwise \(h\)-type, then every sentence
context gives the two fillings the same total \(h\)-value.  Hence their
accepting contexts for \(K(G_0)\) are equal, and \(K(G_0)\) is
\((f,h)\)-tuple-substitutable for every \(f\ge1\).

Conversely, suppose \(L(G_0)\ne\Sigma^*\), and choose \(z\notin L(G_0)\).  Put
\[
  u:=X,
  \qquad v:=Y,
  \qquad E=\blank_1,
  \qquad F=\blank_1\beta z .
\]
Then \(h(u)=h(v)\).  Also \(E[u]=X\in K_{\mathrm{pos}}\setminus Y\Sigma^*\),
because \(X\Sigma^*\cap Y\Sigma^*=\emptyset\), and
\(E[v]=Y\in YL'(G_0)\), because \(\eps\in L'(G_0)\).  Moreover
\(F[u]=X\beta z\) lies in \(K_{\mathrm{pos}}\), since
\(h(X\beta z)=h(Y\beta z)\in h(Y)H_0=S\), and it does not lie in
\(Y\Sigma^*\).  Hence \(F[u]\in K(G_0)\).  On the other hand,
\(F[v]=Y\beta z\) lies in the removed cylinder \(Y\Sigma^*\), and it is not
restored by \(YL'(G_0)\), because \(\beta z\notin L'(G_0)\).  Thus
\(F[v]\notin K(G_0)\).

So \(u\) and \(v\) have the same \(h\)-type and share the accepting context
\(E\), but the context \(F\) separates their distributions.  Therefore the
output language is not \((1,h)\)-tuple-substitutable.  This gives a many-one
reduction from CFG universality to the semantic promise.  The same arity-one
witness refutes \((f,h)\)-tuple-substitutability for every \(f\ge1\).

Finally, failure of \((f,h)\)-tuple-substitutability is semi-decidable by
enumerating finite tuples and contexts and checking the finitely many grammar
membership conditions.  Hence the positive instances are co-r.e.; together
with the universality reduction, they are co-r.e.-complete, equivalently
\(\Pi^0_1\)-complete under standard effective encodings.
\end{proof}

\begin{remark}[Why the promise is semantic]
\label{rem:promise-undecidable}
The preceding theorem sharpens Remark~\ref{rem:semantic-class}.  The learner is
total and consistent on every finite positive sample, but its exactness claim is
necessarily conditional: even at fan-out one, and even for a fixed observation
morphism that is injective on letters, the semantic substitutability promise
cannot be verified from a grammar presentation by any decision procedure.  What
is semi-decidable is only failure, via the finite witnesses displayed in the
proof above.  The theorem is stated for the semantic CFG input problem.  The strictly
\(\eps\)-free working-presentation variant can be handled cleanly in the
marker-letter setting by the same one-marker construction, but the arbitrary-
\(h\) construction above is used here only as a semantic promise result.
\end{remark}

\begin{thebibliography}{99}

\bibitem{Angluin1980}
D.~Angluin.
\newblock Inductive inference of formal languages from positive data.
\newblock {\em Information and Control}, 45(2):117--135, 1980.

\bibitem{Angluin1987}
D.~Angluin.
\newblock Learning regular sets from queries and counterexamples.
\newblock {\em Information and Computation}, 75(2):87--106, 1987.

\bibitem{ClarkEyraud2007}
A.~Clark and R.~Eyraud.
\newblock Polynomial identification in the limit of substitutable context-free
languages.
\newblock {\em Journal of Machine Learning Research}, 8:1725--1745, 2007.

\bibitem{ClarkYoshinaka2014}
A.~Clark and R.~Yoshinaka.
\newblock Distributional learning of parallel multiple context-free grammars.
\newblock {\em Machine Learning}, 96(1--2):5--31, 2014.

\bibitem{ClarkYoshinaka2016}
A.~Clark and R.~Yoshinaka.
\newblock Distributional learning of context-free and multiple context-free
grammars.
\newblock In {\em Topics in Grammatical Inference}, pp.~143--172. Springer,
2016.

\bibitem{delaHiguera1997}
C.~de~la~Higuera.
\newblock Characteristic sets for polynomial grammatical inference.
\newblock {\em Machine Learning}, 27(2):125--138, 1997.

\bibitem{delaHiguera2010}
C.~de~la~Higuera.
\newblock {\em Grammatical Inference: Learning Automata and Grammars}.
\newblock Cambridge University Press, 2010.

\bibitem{DomaratzkiOkhotinShallit2007}
M.~Domaratzki, A.~Okhotin, and J.~Shallit.
\newblock Enumeration of context-free languages and related structures.
\newblock {\em Journal of Automata, Languages and Combinatorics},
12(1--2):79--95, 2007.

\bibitem{Gold1967}
E.~M.~Gold.
\newblock Language identification in the limit.
\newblock {\em Information and Control}, 10(5):447--474, 1967.

\bibitem{Hopcroft1969}
J.~E.~Hopcroft.
\newblock On the equivalence and containment problems for context-free
languages.
\newblock {\em Mathematical Systems Theory}, 3(2):119--124, 1969.
\newblock doi:10.1007/BF01746517.

\bibitem{Kallmeyer2010}
L.~Kallmeyer.
\newblock {\em Parsing Beyond Context-Free Grammars}.
\newblock Springer, Heidelberg, 2010.

\bibitem{Kanazawa1998}
M.~Kanazawa.
\newblock {\em Learnable Classes of Categorial Grammars}.
\newblock CSLI Publications, Stanford, 1998.

\bibitem{kuriyamaCFG}
T.~Kuriyama.
\newblock Distributional learning of context-free languages under fixed finite-monoid typing.
\newblock arXiv:1409.6247v4 [cs.FL], 2026.


\bibitem{LehnerLindorfer2020}
F.~Lehner and C.~Lindorfer.
\newblock Comparing consecutive letter counts in multiple context-free languages.
\newblock arXiv:2002.08236 [cs.FL], 2020.

\bibitem{Shieber1985}
S.~M.~Shieber.
\newblock Evidence against the context-freeness of natural language.
\newblock \emph{Linguistics and Philosophy}, 8(3):333--343, 1985.

\bibitem{Pin1986}
J.-E. Pin.
\newblock {\em Varieties of Formal Languages}.
\newblock North Oxford Academic, London, and Plenum, New York, 1986.

\bibitem{SekiEtAl1991}
H.~Seki, T.~Matsumura, M.~Fujii, and T.~Kasami.
\newblock On multiple context-free grammars.
\newblock {\em Theoretical Computer Science}, 88(2):191--229, 1991.

\bibitem{VijayShankerWeirJoshi1987}
K.~Vijay-Shanker, D.~J. Weir, and A.~K. Joshi.
\newblock Characterizing structural descriptions produced by various grammatical formalisms.
\newblock In {\em Proceedings of the 25th Annual Meeting of the Association for
Computational Linguistics}, pp.~104--111, 1987.

\bibitem{Weir1988}
D.~J. Weir.
\newblock {\em Characterizing Mildly Context-Sensitive Grammar Formalisms}.
\newblock Ph.D. thesis, University of Pennsylvania, 1988.

\bibitem{Yoshinaka2008}
R.~Yoshinaka.
\newblock Identification in the limit of \(k,\ell\)-substitutable context-free
languages.
\newblock In {\em Grammatical Inference: Algorithms and Applications}, LNCS
5278, pp.~266--279. Springer, 2008.

\bibitem{Yoshinaka2009}
R.~Yoshinaka.
\newblock Learning mildly context-sensitive languages with multidimensional
substitutability from positive data.
\newblock In {\em Algorithmic Learning Theory}, LNCS 5809, pp.~278--292.
Springer, 2009.

\bibitem{Yoshinaka2011}
R.~Yoshinaka.
\newblock Efficient learning of multiple context-free languages with
multidimensional substitutability from positive data.
\newblock {\em Theoretical Computer Science}, 412(19):1821--1831, 2011.

\end{thebibliography}

\end{document}